\newpage
\subsection{Banach空间}
\subsubsection{赋范线性空间}
设$X$是数域$\mathbb{K}(\R,\mathbb{C})$上的向量空间,$Z$为$X$的子空间.定义关系
$$
x\sim y\quad \text{当且仅当}\quad (x-y)\in Z
$$
则$\sim$是$X$上的一个等价关系,记所有与$x$等价的元素为$\overline{x}$:
$$
\begin{aligned}
    \overline{x} &= \{y\in X:(x-y)\in Z\}\\
                 &= \{(x-z)\in X:z\in Z\}
\end{aligned}
$$
在商集$X/Z$上定义加法和数乘如下:对任意的$x,y\in X$和$\alpha\in\mathbb{K}$
$$
\overline{x}+\overline{y}=\overline{x+y},\quad \alpha\overline{x}=\overline{\alpha x}
$$
则$X/Z$也成为了数域$\mathbb{K}$上的一个线性空间.\par
\textbf{商范数}$^\clubsuit $:设$(X,\|\cdot\|_{X})$是赋范线性空间,$Z$是$X$的闭子空间,定义映射$\|\cdot\|:X/Z\to\R$为
$$
\|\overline{x}\|=\inf_{y\in\overline{x}}\|y\|_{X}=\inf_{z\in Z}\|x-z\|_{X}
$$
则$\|\cdot\|$是商空间$X/Z$上的范数,称为商范数.\par
\textbf{Proof}:只需要按照范数的定义验证即可:\par
\textbf{Part1.非负性与正定性}.\par

\textbf{压缩映射}$^\clubsuit $:设$(X,d)$为一距离空间.对映射$f:X\to X$,如果存在常数$k$,使得$0<k<1$,且对任何$x,y\in X$均有
$$
d(f(x),f(y))\leqslant kd(x,y)
$$
则称$f$为压缩映射.
\begin{tcolorbox}[colback=KuangNei,colframe=BianKuang,title=\textbf{Banach不动点定理},sharpish corners]
    设$(X,d)$为完备的度量空间,则任何压缩映射$f:X\to X$有且仅有一个不动点$x\in X$.\\
    此外,任意给定点$x_0\in X$,由
    $$
    x_{n+1}=f(x_n),n\geqslant 0
    $$
    定义的序列$(x_n)^{\infty}_{n=0}$当$n\to \infty$时,收敛于$x$,且成立以下估计
    $$
    d(x_n,x)\leqslant \frac{d(f(x_0),x_0)}{1-k}k^n
    $$
\end{tcolorbox}
\textbf{Proof}:\quad \par
我们考虑赋范线性空间中级数的收敛性质.先看下面的例子.\par
\textbf{例}$^\spadesuit$ :设$X=\{x=(\xi_1,\xi_2,\cdots,\xi_n,\cdots):n\text{为任意正整数}\}$.范数$\|x\|=\sup_{k}|\xi_k|$.取
$$
x_k=(0,\cdots,0,\frac{1}{k^2},0,\cdots)\in X
$$
则$\|x_k\|=\frac{1}{k^2}$,因此$\sum_{k=1}^{\infty}\|x_k\|=\frac{\pi^2}{6}$,即$\sum_{k=1}^{\infty}x_k$绝对收敛.但$\sum_{k=1}^{\infty}x_k$在$X$中不收敛,这是因为
$$
\sum_{k=1}^{\infty}x_k=\left(1,\frac{1}{2^2},\cdots,\frac{1}{k^2},\cdots\right)\notin X
$$
这个例子说明,对于一般的赋范线性空间,级数的绝对收敛并不能推出级数本身收敛.幸运的是,我们有下面的定理:\par
\textbf{定理}$^\clubsuit $:设$X$为Banach空间,则从$\sum_{k=1}^{\infty}\|x_k\|$收敛可以推出$\sum_{k=1}^{\infty}x_k$收敛,且满足
$$
\left\|\sum_{k=1}^{\infty}x_k\right\|\leqslant \sum_{k=1}^{\infty}\|x_k\|
$$
反之,在$(X,\|\cdot\|)$中,若任意级数$\sum_{k=1}^{\infty}x_k$绝对收敛总能推出该级数本身收敛,那么$X$是Banach空间.\par
\textbf{Proof}:设$\sum_{k=1}^{\infty}\|x_k\|$收敛,则对任意$\varepsilon>0$,存在$N>0,\forall n>N$,对任何正整数$p$有
$$
\sum_{k=n+1}^{n+p}\|x_k\|<\varepsilon
$$
记$S_n=\sum_{k=1}^n x_k$,于是
$$
\|S_{n+p}-S_n\|=\left\|\sum_{k=n+1}^{n+p}x_k\right\|\leqslant \sum_{k=n+1}^{n+p}\|x_k\|<\varepsilon
$$
这说明$\{S_n\}$是$X$中的Cauchy列,因此在$X$中收敛,即$\sum_{k=1}^{\infty}x_k$收敛.此外,由三角不等式
$$
\left\|\sum_{k=1}^{n}x_k\right\|\leqslant \sum_{k=1}^{n}\|x_k\|
$$
令$n\to\infty$得到$\left\|\sum_{k=1}^{\infty}x_k\right\|\leqslant \sum_{k=1}^{\infty}\|x_k\|$.\par
反之,设$\{S_n\}$是$X$中任意Cauchy列,则对任意$k\in\mathbb{N}$,都存在$n_k>0$,对于$\forall m>n>n_k$有$\|S_m-S_n\|<2^{-k}$.$\forall k\in\mathbb{N}$,都可以选取$n_{k+1}>n_k$.记$x_1=S_{n_1},x_k=S_{n_k}-S_{n_{k-1}}$,于是$\{S_{n_k}\}$是$\sum_{k=1}^{\infty}x_k$的部分和序列.所以
$$
\sum_{k=1}^{\infty}\|x_k\|\leqslant \|x_1\|+\sum_{k=2}^{\infty}2^{-k}=\|x_1\|+1
$$
从而$\sum_{k=1}^{\infty}\|x_k\|$收敛,于是$\sum_{k=1}^{\infty}x_k$收敛,即$\{S_{n_k}\}$收敛,又$\{S_n\}$是$X$中的Cauchy列.所以$\{S_n\}$收敛,于是$X$完备.\par
这个定理的结果给了我们一个证明赋范线性空间完备的方法;这个定理的证明也展示了一个证明空间完备的方法:对于Cauchy列的收敛性不好证明时,我们可以证明其有个子列是收敛的.这个在证明Contor交集定理的逆定理时可以用到.\par
\textbf{Riesz引理}$^\clubsuit $:设$A$赋范线性空间$(X,\|\cdot\|)$的真闭子空间,则对任意$\varepsilon\in (0,1)$,必存在$x_0\in X$且$\|x_0\|=1$使得
$$
d(x_0,A)>\varepsilon
$$
\qquad \textbf{Proof}:由于$A$是$X$的真闭子空间,故存在$x_1\in X\setminus A$,有
$$
d_1=d(x_1,A)=\inf\{\|x_1-x\|:x\in A\}>0
$$
利用下确界性质,对于任意$\delta\in \left(0,\frac{d_1}{\varepsilon}-d_1\right)$,存在$x_2\in A$使得
$$
\|x_1-x_2\|<d_1+\delta<\frac{d_1}{\varepsilon}
$$
记$x_0=\frac{x_1-x_2}{\|x_1-x_2\|}$,从而$\|x_0\|=1$,又$x,x_2\in A$,故$y=\|x_1-x_2\|x+x_2\in A$,且$\|y-x_1\|\geqslant d_1$,于是
$$
\|x-x_0\|=\|x-\frac{x_1-x_2}{\|x_1-x_2\|}\|=\frac{\|y-x_1\|}{\|x_1-x_2\|}\geqslant \frac{d_1}{\|x_1-x_2\|}
$$
从而取下确界可以得到$d(x_0,A)\geqslant \frac{d_1}{\|x_1-x_2\|}>\varepsilon$
\subsubsection{Hilbert空间}
\textbf{定理}$^\clubsuit $:数域$\mathbb{K}$上赋范线性空间$(X,\|\cdot\|)$是内积空间的充要条件是其范数$\|\cdot\|$要满足平行四边形法则
$$
\|x+y\|^2+\|x-y\|^2=2(\|x\|^2+\|y\|^2)
$$
\qquad \textbf{Proof}:必要性.
$$
\begin{aligned}
    \|x+y\|^2 &= (x+y,x+y) = \|x\|^2+\|y\|^2+(x,y)+(y,x)\\
    \|x-y\|^2 &= (x-y,x-y) = \|x\|^2+\|y\|^2-(x,y)-(y,x)
\end{aligned}
$$
将两式相加即可.\par
充分性.考虑构造
$$
(x,y)=
\begin{cases}
    \frac{1}{4}\left(\|x+y\|^2-\|x-y\|^2\right),&\mathbb{K}=\mathbb{R}\\[5pt]
    \frac{1}{4}\left(\|x+y\|^2-\|x-y\|^2+\mathrm{i}\|x+\mathrm{i}y\|^2-\mathrm{i}\|x-\mathrm{i}y\|^2\right),&\mathbb{K}=\mathbb{C}\\
\end{cases}
$$
可以验证,这个就是内积.这两个等式被称作极化恒等式.\par
这个定理表明,由内积导出的范数一定满足极化恒等式,反之,当范数满足平行四边形法则时,由极化恒等式定义的一定是内积.
\begin{tcolorbox}[colback=KuangNei,colframe=BianKuang,title=\textbf{极小化向量定理},sharpish corners]
    设$A$为数域$\mathbb{K}$上内积空间$(X,(\cdot,\cdot))$的非空完备凸子集.对任意给定的$x\in X$,存在唯一的元素$y_0\in A$满足
    $$
    \|x-y_0\|=\inf_{y\in A}\|x-y\|
    $$
    其中范数为由内积导出范数.
\end{tcolorbox}
\textbf{Proof}:如果$x\in A$,定理显然成立.下设$x\notin A$,根据假设$A$非空,因此$\delta=\inf_{y\in A}\|x-y\|$存在且$\delta\geqslant 0$.若$\delta= 0$,则$x\in\overline{A}$,而$A$是闭集,必有$x\in A$,这就矛盾了.从而$\delta>0$.\par
由确界的性质,存在$y_n\in A,n\geqslant 0$使得
$$
\|x-y_n\|\to \delta=\inf_{y\in A}\|x-y\|,\quad n\to\infty
$$
由平行四边形法则,对任何的$m,n\geqslant 0$,
$$
\begin{aligned}
    \|y_m-y_n\|^2 &=\|(x-y_m)-(x-y_n)\|^2\\
                  &= 2\|x-y_m\|^2+2\|x-y_n\|^2-\|2x-(y_m+y_n)\|^2\\
                  &= 2\|x-y_m\|^2+2\|x-y_n\|^2-4\left\|x-\frac{y_m+y_n}{2}\right\|^2
\end{aligned}
$$
而集合$A$是凸集,从而$\frac{y_m+y_n}{2}\in A$,因此
$$
\left\|x-\frac{y_m+y_n}{2}\right\|^2\geqslant \delta ^2
$$
因此,对任何的$m,n\geqslant 0$,
$$
0\leqslant \|y_m-y_n\|^2\leqslant 2\|x-y_m\|^2+2\|x-y_n\|^2-4\delta ^2
$$
故而$\{y_n\}_{n\geqslant 0}$是Cauchy列.由假设,$A$是完备的,存在$y_0\in A$使得$y_n\to y_0$.利用范数的连续性知
$$
\|x-y_0\|=\lim_{n\to\infty}\|x-y_n\|=\delta=\inf_{y\in A}\|x-y\|
$$
若还有$z_0\in A$使得$\|x-z_0\|=\delta$,则
$$
\begin{aligned}
    \|y_0-z_0\|^2 &= 2\left(\|y_0-x\|^2+\|x-z_0\|^2\right)-\frac{1}{4}\left\|x-\frac{1}{2}(y_0+z_0)\right\|^2\\
                  &\leqslant 4\delta^2-4\delta^2=0
\end{aligned}
$$
由此可以得出$z_0=y_0$.\par
\begin{tcolorbox}[colback=KuangNei,colframe=BianKuang,title=\textbf{正交分解定理},sharpish corners]
    设$A$为数域$\mathbb{K}$上内积空间$(X,(\cdot,\cdot))$的非空完备子空间.对任意给定的$x\in X$,存在唯一的正交分解
    $$
    x=y_0+z,\quad y_0\in A,z\in A^{\perp }
    $$
\end{tcolorbox}
\textbf{Proof}:因为子空间都是凸集,所以$\forall x\in X$,存在唯一的$y_0\in A$使得
$$
\|x-y_0\|=\inf\{\|x-y\|:y\in A\}
$$
令$z=x-y_0$,下证$z\in A^{\perp }$,即$\forall y\in A,(z,y)=0$.由于$(z,0)=0$,故不妨设$y\neq 0$.令$(y,z)=\alpha$,则$\forall \lambda\in\mathbb{K}$有
$$
\|z-\lambda y\|^2=(z-\lambda y,z-\lambda y)=\|z\|^2+\lambda\overline{\lambda}\|y\|^2-\lambda\alpha-\overline{\lambda}\overline{\alpha}
$$
取$\lambda=\frac{\overline{\alpha}}{\|y\|^2}$,代入上式可得
$$
\|z-\lambda y\|^2=\|z\|^2-\frac{|\alpha|^2}{\|y\|^2}
$$
另一方面,$y_0+\lambda y\in A$,因此
$$
\|z\|=\|x-y_0\|\leqslant \|x-(y_0\lambda y)\|=\|z-\lambda y\|
$$
于是$\|z\|^2-\frac{|\alpha|^2}{\|y\|^2}\geqslant \|z\|^2$,从而$\|\alpha\|=0$,即$(z,y)=0$.\par
\textbf{定义}$^\clubsuit $:设$A$是内积空间$X$的子空间,$x\in X$,若存在$y\in A,z\in A^{\perp }$,使得$x=y+z$,则称$y$是$x$在$A$上的正交投影,记为$y=Px$.\par
需要注意的是,对于一般内积空间$X$的任意子空间$A$,$X$中向量$x$在$A$上的投影并不一定存在.从极小化向量定理可以看出,$x\in X$在$A$中的最佳逼近元$y_0$就是$x$在$A$上的投影,而且当$A$是内积空间$X$的完备子空间时,$x$在$A$中的投影必唯一存在.\par
下面的讨论中,我们只考虑$\mathbb{K}=\R$的情形.\par
\textbf{推论1}$^\clubsuit $:$y\in A$,$y=Px$的充分必要条件是对任意$z\in A$
$$
(y-x,z-y)\geqslant 0
$$
\qquad\textbf{Proof}:\textbf{必要性}.\par
如果$x\in A$,此时$Px-x=0$,欲证的不等式显然成立.下设$x\notin A$,给定$y\in A$,对于$0\leqslant\theta\leqslant 1,y+\theta(y-y)\in A$.由$y=Px$的定义可得,对于$0\leqslant\theta\leqslant 1$
$$
\begin{aligned}
    \|x-y\|^2& \leqslant \|x-(y+\theta(y-y))\|^2\\
             &= \|x-y\|^2-2\theta(x-y,y-y)+\theta ^2\|y-y\|^2
\end{aligned}
$$
因此,对于$0\leqslant\theta\leqslant 1$,
$$
0\leqslant 2\theta(y-x,y-y)+\theta ^2\|y-y\|^2
$$
这是一个关于$\theta$的二次函数,可以得到$(y-x,y-y)\geqslant 0$.\par
\textbf{充分性}.\par
设$y\in A$使得对任何$y\in A,(y-x,y-y)\geqslant 0$,于是对任何$y\in A$,
$$
\begin{aligned}
    \|x-y\|^2 &= \|x-y+y-y\|^2=\|x-y\|^2+2(y-x,y-y)+\|y-y\|^2\\
              &\geqslant \|x-y\|^2
\end{aligned}
$$
由此得到$y=Px$.\par
\textbf{推论2}$^\clubsuit $:投影定理中,我们定义了一个映射$P:X\to A$,它满足对所有的$x_1,x_2\in X$
$$
\|Px_1-Px_2\|\leqslant\|x_1-x_2\|
$$
即$P$是Lipschity连续映射.\par
\textbf{Proof}:根据推论1,由$Px_2\in A$得到
$$
(Px_1-x_1,Px_2-Px_1)\geqslant 0
$$
因为$Px_1\in A$,所以
$$
(x_2-Px_2,Px_2-Px_1)\geqslant 0
$$
从而
$$
(Px_1-Px_2,Px_2-Px_1)+(x_2-x_1,Px_2-Px_1)\geqslant 0
$$
由Cauchy-Schwary不等式可得
$$
\|Px_1-Px_2\|^2\leqslant (x_2-x_1,Px_2-Px_1)\leqslant \|x_2-x_1\|\|Px_2-Px_1\|
$$
这就证明了不等式.\par
\textbf{推论3}$^\clubsuit $:设子集$A$是$X$的完备子空间,则对所有的$y\in A$,
$$
(Px-x,y)=0
$$
反之,若元素$y\in A$满足对所有的$y\in A$,
$$
(y-x,y)=0
$$
那么$y=Px$.\par
\textbf{Proof}:设$x\notin A$,给定$y\in A$,对所有的$\theta\in\R,Px+\theta y\in A$,由推论1可得对所有的$\theta\in\R$
$$
(Px-x-Px+\theta y-Px)=\theta(Px-x,y)\geqslant 0
$$
因此,$(Px-x,y)=0$.\par
反之,设$y\in A$满足对任何$y\in A$,有$(y-x,y)=0$,特别的有$(y-x,x)=0$,因此对任何$y\in A$,有
$$
(y-x,y-y)=0\geqslant 0
$$
由推论1可得$y=Px$.\par
\begin{tcolorbox}[colback=KuangNei,colframe=BianKuang,title=\textbf{Gram-Schmidt方法},sharpish corners]
    设$\{x_k\}$是内积空间$X$中任一线性无关集,则由$\{x_k\}$必可作出一个标准正交系$\{e_k\}$,使得对任意$n$有
    $$
    \operatorname{span}\{e_1,\cdots,e_n\}=\operatorname{span}\{x_1,\cdots,x_n\}
    $$
\end{tcolorbox}
\textbf{Proof}:设$e_1=\frac{x_1}{\|x_1\|}$,则$\|e_1\|=1$.令$y_1=x_1,y_2=x_2-(x_2,e_1)e_1$,则$y_2\neq 0$.于是可令$e_2=\frac{y_2}{\|y_2\|}$,则$\|e_2\|=1$,且$(y_2,e_1)=(x_2,e_1)-(x_2,e_1)=0$,即有$(e_1,e_2)=0$且$\operatorname{span}\{e_1,e_2\}=\operatorname{span}\{x_1,x_2\}$.\par
按上面的方法作出$e_1,\cdots,e_{n-1}$.再令
$$
y_n=x_n-\sum_{k=1}^{n-1}(x_n,e_k)e_k
$$
由$x_1,\cdots,x_n$线性无关,那么$y_n\neq 0$,于是令$e_n=\frac{y_n}{\|y_n\|}$,则$\|e_n\|=1$,且$(e_n,e_k)=0,k=1,\cdots,n-1$.\par
一方面
$$
x_n=\sum_{k=1}^{n-1}(x_n,e_k)e_k+y_n=\sum_{k=1}^{n-1}(x_n,e_k)e_k+\|y_n\|e_n=\sum_{k=1}^{n}c_ke_k
$$
其中$c_k=(x_n,e_k),k=1,\cdots,n-1,c_n=\|y_n\|$.另一方面
$$
e_n=\frac{y_n}{\|y_n\|}=\frac{1}{\|y_n\|}\left(x_n-\sum_{k=1}^{n-1}(x_n,e_k)e_k\right)
$$
由归纳法可以可以求得$e_n=\sum_{k=1}^{n}\alpha_k x_k$.这说明$\{x_1,\cdots,x_n\}$和$\{e_1,\cdots,e_n\}$可以互相线性表示,从而$\operatorname{span}\{e_1,\cdots,e_n\}=\operatorname{span}\{x_1,\cdots,x_n\}$.\par
我们现在的目光转向内积空间的Fourier展开.\par
\textbf{最佳逼近定理}$^\clubsuit $:设$\{e_k\}$是内积空间$X$中的标准正交系,$x\in X,c_k=(x,e_k),A_n=\operatorname{span}\{e_1,\cdots,e_n\}$,则Fourier级数的部分和$S_n=\sum_{k=1}^{n}c_ke_k$是$x$在$A_n$中的最佳逼近元,即
$$
\|x-S_n\|=\min\left\{\|x-y_n\|:y_n=\sum_{k=1}^{n}\alpha_ke_k\in A_n,\alpha\in \mathbb{K}\right\}
$$
\qquad \textbf{Proof}:对于任意$j,1\leqslant j\leqslant n$,有
$$
(x-S_n,e_j)=(x,e_j)-\sum_{k=1}^{n}c_k(e_k,e_j)=(x,e_j-c_j)=0
$$
即$(x-S_n,e_j)=0$.从而$\forall y_n=\sum_{k=1}^{n}\alpha_ke_k\in A_n,(x-S_n,y_n)=0$.利用勾股定理
$$
\begin{aligned}
\|x-y_n\|^2 &= \|(x-S_n)+(S_n-y_n)\|^2\\
            &= \|x-S_n\|^2+\left\|\sum_{k=1}^{n}(c_k-\alpha_k)e_k\right\|^2\\
            &= \|x-S_n\|^2+\sum_{k=1}^{n}|c_k-\alpha_k|^2
\end{aligned}
$$
所以当且仅当$c_k=\alpha_k$时,$\|x-y_n\|$达到最小值,且最小值为$\|x-S_n\|$.\par
\begin{tcolorbox}[colback=KuangNei,colframe=BianKuang,title=\textbf{Bessel不等式},sharpish corners]
    设$\{e_k\}$是内积空间$X$中的标准正交系,则$\forall x\in X,c_k=(x,e_k)$成立Bessel不等式
    $$
    \sum_{k=1}^{\infty}|c_k|^2\leqslant \|x\|^2
    $$
\end{tcolorbox}
\textbf{Proof}:$\forall n,S_n=\sum_{k=1}^{n}c_ke_k$,则
$$
\begin{aligned}
    \|x\|^2=\|x-S_n+S_n\|^2 &= \|x-S_n\|^2+\|S_n\|^2\\
                            &= \left\|x-\sum_{k=1}^{n}c_ke_k\right\|^2+\sum_{k=1}^{n}|c_k|^2
\end{aligned}
$$
即$\|x\|^2-\sum_{k=1}^{n}|c_k|^2=\left\|x-\sum_{k=1}^{n}c_ke_k\right\|^2\geqslant 0$.从而$\sum_{k=1}^{n}|c_k|^2\leqslant \|x\|^2$.两边令$n\to\infty$就得到Bessel不等式.\par
\begin{tcolorbox}[colback=KuangNei,colframe=BianKuang,title=\textbf{Parseval等式},sharpish corners]
    设$A=\{e_k\}$是内积空间$X$中的标准正交系,则$\forall x\in X,c_k=(x,e_k)$为$x$的Fourier系数,则下面的命题等价:\\
    (1).$\overline{\operatorname{span}A}=X$,即$A=\{e_k\}$是$X$中的标准正交基;\\
    (2).成立Parseval等式:$\sum_{k=1}^{\infty}|c_k|^2=\|x\|^2$;\\
    (3).$x=\sum_{k=1}^{\infty}c_ke_k$.
\end{tcolorbox}
\textbf{Proof}:我们采用"(1)$\Rightarrow $(2) $\Rightarrow$ (3) $\Rightarrow$ (1)".\par
\textbf{Part1.(1)$\Rightarrow $(2)}\par
设$\overline{\operatorname{span}A}=X$,则对任意$x\in X$,存在$y_n=\sum_{k=1}^{n}\alpha_ke_k$使得
$$
\|x-y_n\|\to 0\quad (n\to\infty)
$$
又(最佳逼近定理)$\|x-S_n\|\leqslant \|x-y_n\|$.从而$\|x-S_n\|\to 0$.于是由Bessel不等式的证明过程得
$$
\|x\|^2-\sum_{k=1}^{n}|c_k|^2=\|x-S_n\|^2\to 0,\quad (n\to\infty)
$$
\qquad \textbf{Part2.(2)$\Rightarrow $(3)}\par
设$\sum_{k=1}^{\infty}|c_k|^2=\|x\|^2$,即$\sum_{k=1}^{n}|c_k|^2-\|x\|^2\to 0(n\to \infty)$.又
$$
\|x-S_n\|^2=\|x\|^2-\sum_{k=1}^{n}|c_k|^2\to 0\quad (n\to\infty)
$$
也即有$x=\sum_{k=1}^{\infty}c_ke_k$.\par
\textbf{Part3.(3)$\Rightarrow $(1)}\par
设$\forall x\in X$,成立$\|x-S_n\|\to 0(n\to\infty)$.由于$S_n=\sum_{k=1}^{n}c_ke_k\in \operatorname{span}A$,所以$x\in \overline{\operatorname{span}A}$,从而有$X\subset \overline{\operatorname{span}A}$.显然有$\overline{\operatorname{span}A}\subset X$,因此$\overline{\operatorname{span}A}=X$.\par
\begin{tcolorbox}[colback=KuangNei,colframe=BianKuang,title=\textbf{Riesz表示定理},sharpish corners]
    设$H$为数域$\mathbb{K}$上的Hilbert空间,其上的内积记为$(\cdot,\cdot)$.对于任意给定的连续线性泛函$f:H\to\mathbb{K}$,存在唯一的$y\in H$,使得对$\forall x\in H$
    $$
    \langle f,x\rangle=(x,y)
    $$
    而且$\|f\|=\|y\|$
\end{tcolorbox}
\textbf{Proof}:因为$f\in H'$,于是$f$的零空间$N(f)=\{x\in H:f(x)=0\}$是$H$的闭子空间.若$N(f)=H$,则$f=0$,这时只要取$y=0$即可;下设$N(f)\neq H$,由正交分解定理,存在$z\in N(f)^{\perp },z\neq 0$,所以$z\notin N(f)$(这是因为$N(f)\cap N(f)^{\perp}=\{0\}$),即$f(z)\neq 0$.\par
令$A=\{u=zf(x)-xf(z):x\in H\}$,对于$u\in A$,$f(u)=f(x)f(z)-f(z)f(u)=0$,所以$A\subset N(f)$,从而$z\in N(f)^{\perp}\subset A^{\perp}$,于是
$$
(u,z)=0,\quad \forall x\in H
$$
将$u=zf(x)-xf(z)$代入上式可得$f(x)=\frac{f(z)}{\|z\|^2}(x,z)$,令$y=\frac{\overline{f(z)}}{\|z\|^2}$即得$f(x)=(x,y)$.\par
由Schwarz不等式
$$
\|f\|=\sup_{\|x\|=1}|f(x)|=\sup_{\|x\|=1}|(x,y)|\leqslant \|y\|
$$
另一方面,$\|y\|^2=|(y,y)|=f(y)\leqslant \|f\|\|y\|$,由$y\neq 0$,所以$\|f\|\geqslant \|y\|$.综上我们得到了
$$
\|f\|= \|y\|
$$
\qquad 最后,我们证明定理中的$y$是唯一的.若还存在$y'\in H$使得$f(x)=(x,y'),\forall x\in H$,则$(x,y)=(x,y')$,即$(x,y-y')=0$,取$x=y-y'$得到
$$
\|y-y'\|^2=(y-y',y-y')=0
$$
因此$y=y'$.
\newpage
\subsection{线性算子}
定义不再赘述,我们只给出一些重要的定理和性质.\par
\textbf{Contor交集定理}$^\clubsuit $:设$(X,d)$是完备度量空间,$\{A_n\}_{n=1}^{\infty}$为$X$中满足下面性质的非空闭子集$A_n$组成的序列:
$$
A_0\supset A_1\supset A_2\supset \cdots\text{且当}n\to\infty\text{时},\operatorname{diam}A_n\to 0
$$
则存在$x\in X$使得
$$
\bigcap_{n=1}^{\infty}A_n=\{x\}
$$
\qquad \textbf{Proof}:对每个$n\geqslant 0$,取一个元素$x_n\in A_n$(这是因为每个$A_n$都非空).由于对所有的$m\geqslant n$都有$A_m\subset A_n$,这就表明了
$$
d(x_m,x_n)\leqslant \operatorname{diam}A_n,\forall m\geqslant n
$$
又当$n\to\infty$时,有$\operatorname{diam}A_n\to 0$,所以$\{x_n\}_{n=0}^{\infty}$是Cauchy列.设$x=\lim\limits_{n\to\infty}x_n$.\par
给定任一$n\geqslant 0$,对所有的$m\geqslant n$,都有$x_m\in A_n$.又因为$A_n$是闭集,所以$x=\lim\limits_{m\to\infty}x_m\in A_n$.这就说明了$x\in \bigcap_{n=1}^{\infty}A_n$.\par
假设还有一点$y\in \bigcap_{n=1}^{\infty}A_n$且$x\neq y$,由于当$n\to\infty$时,$\operatorname{diam}A_n\to 0$,而$d(x,y)>0$,所以存在$n_0>0$使得$\operatorname{diam}A_{n_0}<d(x,y)$.但是因为$x,y$都属于$A_{n_0}$,所以$d(x,y)\leqslant\operatorname{diam}A_{n_0}$,这就推出矛盾,只能有$x=y$.因此$\bigcap_{n=1}^{\infty}A_n=\{x\}$.\par
\begin{tcolorbox}[colback=KuangNei,colframe=BianKuang,title=\textbf{Baire定理},sharpish corners]
    设$X$是完备度量空间,则下述两个等价性质成立:\\
    (a)\, 设$(F_n)^{\infty}_{n=0}$是$X$中的闭子集序列,且对所有的$n\geqslant 0,\operatorname{int}\,F_n=\emptyset$,则
    $$
    \operatorname{int}\,\left(\bigcup^{\infty}_{n=0}F_n\right)=\emptyset
    $$
    (b)\, 设$(O_n)^{\infty}_{n=0}$是$X$中的开子集序列,且对所有的$n\geqslant 0,\overline{O}_n=X$,则
    $$
    \overline{\bigcap^{\infty}_{n=0}O_n}=X
    $$
\end{tcolorbox}
\textbf{Proof}:定理的证明分成两个部分.\par
\textbf{Part1.先证明(a)和(b)等价.设(a)成立}.\par
注意到,对任意的$A\subset X,\overline{A}=X-\operatorname{int}(X-A)$,这说明
$$
\overline{A}=X\text{当且仅当}\operatorname{int}(X-A)=\emptyset
$$
设开集$O_n\subset X,n\geqslant 0$,使得对所有的$n\geqslant 0,\overline{O}_n=X$.设$F_n:=X-O_n$,则闭集$F_n$就满足对所有的$n\geqslant 0,\operatorname{int}F_n=\emptyset$.所以由(a)得$\operatorname{int}\,\left(\bigcup^{\infty}_{n=0}F_n\right)=\emptyset$.再由de Morgan定律:
$$
\bigcup^{\infty}_{n=0}F_n=\bigcup^{\infty}_{n=0}(X-O_n)=X-\bigcap^{\infty}_{n=0}O_n
$$
因此$\operatorname{int}\left(X-\bigcap^{\infty}_{n=0}O_n\right)=\emptyset$.这就有$\overline{\bigcap^{\infty}_{n=0}O_n}=X$.故(b)成立.\par
注意到对任意的$B\subset X,\overline{X-B}=X-\operatorname{int}\,B$,这说明
$$
\operatorname{int}\,B=\emptyset\text{当且仅当}\overline{X-B}=X
$$
现给定闭集$F_n\subset X,n\geqslant 0$,使得对所有的$n\geqslant 0,\operatorname{int}\,F_n=\emptyset$.设$O_n:=X-F_n$,则开集$O_n$就满足对所有的$n\geqslant 0,\overline{O_n}=X$.所以由(b)得$\overline{\bigcap^{\infty}_{n=0}O_n}=X$.再由de Morgan定律:
$$
\bigcap^{\infty}_{n=0}O_n=\bigcap^{\infty}_{n=0}(X-F_n)=X-\bigcup^{\infty}_{n=0}F_n
$$
因此$\overline{X-\bigcup^{\infty}_{n=0}F_n}=X$,这就有$\operatorname{int}\,\left(\bigcup^{\infty}_{n=0}F_n\right)=\emptyset$.故(a)成立.\par
\textbf{Part2.现在证明(a)成立}.\par
注意到$X$的任一子集$A$有非空内部的充分必要条件是存在一个包含于$A$中的非空开子集$O\subset X$或者等价地$O\cap (X-A)=\emptyset$.于是其逆否命题就是
$$
\operatorname{int}\,A=\emptyset\text{的充分必要条件是对所有的非空开子集}O\subset X,O\cap (X-A)\neq \emptyset
$$
\qquad 设$X$是完备度量空间,$F_n(n\geqslant 0)$是$X$的闭子集使得对所有的$n\geqslant 0$有$\operatorname{int}\, F_n=\emptyset$.我们希望证明的是$\operatorname{int}\left(\bigcup^{\infty}_{n=0}F_n\right)=\emptyset$,或者等价地证明对任意给定的非空开集$O\subset X$,均有
$$
O\cap\left(X-\bigcup^{\infty}_{n=0}F_n\right)\neq \emptyset
$$
\qquad 给定一个非空开集$O\subset X$,设$O_0=O$.由于$\operatorname{int}\,F_0=\emptyset$且$O_0$是开集,故$O_0\cap (X-F_0)$是$X$的一个非空开子集,所以存在非空开子集$O_1\subset X$(存在性根据开集定义即可得到)使得
$$
\overline{O}_1\subset O_0\cap(X-F_0)\text{且}\operatorname{diam}\overline{O}_1<1
$$
因为$\operatorname{int}\,F_1=\emptyset$且$O_1$是开集,故$O_1\cap (X-F_1)$是$X$的非空开子集.所以存在非空开子集$O_2\subset X$使得
$$
\overline{O}_2\subset O_1\cap(X-F_1)\text{且}\operatorname{diam}\overline{O}_2<\frac{1}{2}
$$
用这样的方式继续下去,存在非空开子集$O_{n+1}\subset X,n\geqslant 0$,使得
$$
\overline{O}_{n+1}\subset O_n\cap(X-F_n)\text{且}\operatorname{diam}\overline{O}_{n+1}<\frac{1}{n+1},n\geqslant 0
$$
\qquad 非空闭子集满足Contor交集定理的假设,所以存在$x\in X$使得$\{x\}=\bigcap_{n=0}^{\infty}\overline{O}_n$.\par
显然$x\in \overline{O}_1\subset O_0=O$.另一方面,由于$x\in \overline{O}_{n+1}$且对所有的$\geqslant 0$成立$\overline{O}_{n+1}\subset X-F_n$,这表明$x\in \bigcap_{n=1}^{\infty}(X-F_n)=X-\bigcup_{n=1}^{\infty}F_n$.我们有
$$
x\in O\cap\left(X-\bigcup^{\infty}_{n=0}F_n\right)
$$
这就证明了对任意给定的非空开集$O\subset X$,都有$O\cap\left(X-\bigcup^{\infty}_{n=0}F_n\right)\neq \emptyset$.\par
在应用过程中,我们常采用下面形式的Baire定理:\par
\textbf{推论}$^\clubsuit $:设$X$是度量空间,令$(F_n)^{\infty}_{n=0}$是$X$中的闭子集序列使得$X=\bigcup^{\infty}_{n=0}F_n$.\par
\quad (a)如对所有的$n\geqslant 0,\operatorname{int}\,F_n=\emptyset$,则$X$是不完备的.\par
\quad (b)如果$X$是完备的,则存在$n_0\geqslant 0$,使得$\operatorname{int} F_{n_0}\neq \emptyset$.\par
\begin{tcolorbox}[colback=KuangNei,colframe=BianKuang,title=\textbf{Banach-Steinhaus定理(一致有界性定理)},sharpish corners]
    设$X$是Banach空间,$Y$是赋范向量空间,而$(T_i)_{i\in I}$为映射$T_i\in\mathcal{L}(X;Y)$构成的算子族,满足
    $$
    \sup\limits_{i\in I}\|T_ix\|_{Y}<\infty,\forall x\in X
    $$
    则
    $$
    \sup\limits_{i\in I}\|T_i\|_{\mathcal{L}(X;Y)}<\infty
    $$
\end{tcolorbox}
\textbf{Proof}:对每个整数$n\geqslant 0$,定义集合
$$
F_n=\{x\in X,\sup\limits_{i\in I}\|T_ix\|_{Y}\leqslant n\}
$$
给定任意的$x\in X$,由假设$\sup\limits_{i\in I}\|T_ix\|_{Y}<\infty$,故存在一个整数$n(x)\geqslant 0$使得$\sup\limits_{i\in I}\|T_ix\|_{Y}\leqslant n(x)$.这表明$x\in F_{n(x)}$.所以
$$
X=\bigcup_{n=0}^{\infty}F_n
$$
根据定义,$x\in F_n$的充分必要条件是对所有的$i\in I$有$x\in \{z\in X,\|T_iz\|_{Y}\leqslant n\}$.所以,集合$F_n$又可以表示为
$$
F_n=\bigcap_{i\in I}\{z\in X,\|T_iz\|_{Y}\leqslant n\}
$$
由于每个线性算子$T_i$都是连续的,故集合$\{z\in X,\|T_iz\|_{Y}\leqslant n\}$是闭集,从而$F_n$是$X$中的闭集.\par
因为$X$是完备的,根据Baire定理推论(b),存在一个整数$n_0\geqslant 0$使得$\operatorname{int}\,F_{n_0}\neq\emptyset$.所以存在$x_0\in F_{n_0}$及$r>0$使得$\overline{B(x_0,r)}\subset F_{n_0}$.再由$F_{n_0}$的定义,这意味着
$$
\|T_iz\|_{Y}\leqslant n_0,\forall z\in \overline{B(x_0,r)},\forall i\in I
$$
对任意非零向量$x\in X$,都可以写成
$$
x=\frac{\|x\|}{r}(z-x_0),\text{其中}z=\left(x_0+\frac{r}{\|x\|}x\right)\in\overline{B(x_0,r)}
$$
这就有
$$
\begin{aligned}
    \|T_ix\|_{Y} &\leqslant \frac{\|x\|}{r}(\|T_iz\|_{Y}+\|T_ix_0\|_{Y})\\
                 &\leqslant \frac{1}{r}(n_0+\|T_ix_0\|_{Y})\|x\|\\
                 &\leqslant \frac{1}{r}(n_0+\sup_{i\in I}\|T_ix_0\|_{Y})\|x\|,\forall i\in I,x\in X
\end{aligned}
$$
由假设$\sup_{i\in I}\|T_ix\|_{Y}<\infty$,所以我们有
$$
\sup_{i\in I}\|T_i\|_{Y}\leqslant \frac{1}{r}(n_0+\sup_{i\in I}\|T_ix_0\|_{Y})<\infty
$$
\begin{tcolorbox}[colback=KuangNei,colframe=BianKuang,title=\textbf{Banach开映射定理},sharpish corners]
    设$X$及$Y$是Banach空间,而$T\in\mathcal{L}(X;Y)$是满射.则$X$的任一开子集$U$在映射$T$下的像$T(U)$是$Y$的开子集.
\end{tcolorbox}
\textbf{Proof}:为了不在证明中混淆,我们将$X$中球心在$X$原点的开球,以及$Y$中的开球分别用下面的记号表示:
$$
B_r=\{x\in X:\|x\|<r\},\quad B(y,s)=\{\tilde{y}\in Y:\|\tilde{y}-y\|<s\}
$$
给定一个向量空间$Z$,向量$z_0\in Z$,数$\alpha$及一个子集$A\subset Z$,设
$$
\{z_0\}+A=\{z_0+z:z\in A\},\quad \alpha A=\{\alpha z:z\in A\}
$$
因此若$z_0\in A$,那么有$\{z_0\}+A\subset 2A$.\par
证明分成四步.\par
\textbf{Part1:集合$\overline{T(B_1)}$包含一个开球}.\par
给定$y\in Y$,因为$T$是满射,故存在$x\in X$,使得$y=Tx$.由于$x\in B_n$一定对某个整数$n\geqslant 1$成立,所以$Y=\bigcup_{n=1}^{\infty}T(B_n)$,从而必有
$$
Y=\bigcup_{n=1}^{\infty}\overline{T(B_n)}
$$
由于$Y$完备,由Baire定理的推论可知,存在整数$n_0\geqslant 1$使$\operatorname{int}\overline{T(B_{n_0})}\neq \emptyset$.由$T$的线性性质知$\overline{T(B_1)}=\frac{1}{n_0}\overline{T(B_{n_0})}$.
故$\operatorname{int}\overline{T(B_1)}\neq\emptyset$.这说明$\overline{T(B_1)}$包含一个开球.\par
\textbf{Part2.集合$\overline{T(B_1)}$包含一个球心在$Y$原点的开球}.\par
根据Part1,存在$y\in Y$及$s>0$使得$B(y,2s)\subset\overline{T(B_1)}$.根据定义可以得到
$$
B(0,2s)=\{-y\}+B(y,2s)\subset \{-y\}+\overline{T(B_1)}
$$
因为$y\in \overline{T(B_1)}$,所以存在序列$\{y_n\}\subset T(B_1)$使得$\lim\limits_{n\to\infty}y_n=y$.我们需要明白$T(B_1)$的定义是
$$
T(B_1)=\{T(x):x\in B_1\}
$$
因此$y_n\in T(B_1)$的意思根据定义就是,存在$x_n\in B_1$使得$y_n=T(x_n)$.显然$-x_n\in B_1$.于是$T(-x_n)=-T(x_n)=-y_n$.于是我们找到了序列$\{-y_n\}\subset T(B_1)$,且$\lim\limits_{n\to\infty}(-y_n)=-y$.这说明$-y\in\overline{T(B_1)}$.从而得到
$$
B(0,2s)=\{-y\}+B(y,2s)\subset \{-y\}+\overline{T(B_1)}\subset 2\overline{T(B_1)}
$$
再根据$T$是线性的,就有$B(0,s)\subset \overline{T(B_1)}$.\par
\textbf{Part3.集合$T(B_1)$包含一个球心在$Y$原点的开球}.\par
为此,我们将证明$B\left(0,\frac{s}{2}\right)\subset T(B_1)$,这里的$s>0$和Part2一致.根据定义,这意味着,给定任意的$y\in B\left(0,\frac{s}{2}\right)$,我们要找到$x\in B_1$,使得$y=Tx$.现设$y\in B\left(0,\frac{s}{2}\right)$给定.\par
根据Part2,有$y\in B\left(0,\frac{s}{2}\right)\subset\overline{T(B_{1/2})}$,故存在$x_1\in B_{1/2}$使得$\|y-Ax_1\|<\frac{s}{2^2}$;再根据Part2,有
$$
(y-Tx_1)\in B\left(0,\frac{s}{2^2}\right)\subset \overline{T(B_{1/2^2})}
$$
故存在$x_2\in B_{1/2^2}$使得$\|y-Tx_1-Tx_2\|<\frac{s}{2^3}$.按这种方式,我们将构造出序列$x_n\in X$,它满足
$$
x_n\in B_{1/2^n},\text{并且}\left\|y-T\left(\sum_{k=1}^nx_k\right)\right\|<\frac{s}{2^{n+1}},\forall n\geqslant 1
$$
由于对每个$n\geqslant 1,\|x_n\|<\frac{1}{2^n}$,所以级数$\sum_{n=1}^{\infty}x_n$是绝对收敛的.而空间$X$完备,所以存在$x\in X$使得
$$
\lim\limits_{n\to\infty}\sum_{k=1}^nx_k=x,\text{并且}\|x\|\leqslant \sum_{k=1}^{\infty}\|x_k\|<\frac{1}{2}+\frac{1}{2^2}+\cdots+\frac{1}{2^{n+1}}+\cdots=1
$$
故$x\in B_1$,又因为$T$是连续的,所以
$$
y=\lim\limits_{n\to\infty}T\left(\sum_{k=1}^nx_k\right)=Tx
$$
\qquad \textbf{Part4.算子$T$是开的.}\par
给定$X$的任一开子集$U$及任意点$y\in T(U)$,我们必须找到$\delta>0$使得$B(y,\delta)\subset T(U)$.现设$x\in U$使$y=Tx$.\par
因为$U$是开的,所以存在$r>0$,使得$B_r(x)=\{x\}+B_r\subset U$.根据Part3,存在$\delta>0$使得$B(0,\delta)\subset T(B_r)$.因此
$$
B(y,\delta)=\{y\}+B(0,\delta)\subset \{y\}+T(B_r)=T(\{x\}+B_r)=T(B_r(x))\subset T(U)
$$
综上我们证明了算子$T$是开的.\par
由Banach开映射定理容易得到下面的推论,它常被用作判定线性算子逆的连续性的充分条件.\par
\textbf{逆算子定理}$^\clubsuit $:设$X$及$Y$是Banach空间(它们都是数域$\mathbb{K}$上的),而$T\in\mathcal{L}(X;Y)$是双射.则$T^{-1}\in\mathcal{L}(Y;X)$.\par
\textbf{Proof}:因为$T$是双射,逆映射$T^{-1}$的存在性自然成立.接下来我们需要证明两件事情:$T^{-1}$是线性的和连续的.\par
\textbf{Part1.$T^{-1}$是线性的}.\par
任意给定两个向量$y_1,y_2\in Y$,存在唯一确定的向量$x_1,x_2\in X$,使得$y_1=Tx_1,y_2=Tx_2$.对任意的$\alpha,\beta\in\mathbb{K}$,
$$
\begin{aligned}
    T^{-1}(\alpha y_1+\beta y_2) &= T^{-1}(\alpha Tx_1+\beta Tx_2)\\
                                 &= T^{-1}(T(\alpha x_1+\beta x_2))=\alpha x_1+\beta x_2=\alpha T^{-1}y_1+\beta T^{-1}y_2
\end{aligned}
$$
\qquad \textbf{Part2.$T^{-1}$是连续的}.\par
我们只需要证明$T^{-1}$在原点是连续的即可.根据定义,只要证明给定任意的开球$B_r\subset X$,存在一个开球$B(0,\delta)\subset Y$,使得
$$
T^{-1}(B(0,\delta))\subset B_r
$$
由于$T$是双射,上述包含关系等价于$B(0,\delta)\subset T(B_r)$.而在Banach开映射定理的证明中,我们已知存在这样的$\delta>0$.\par
Banach开映射定理的下一个推论给出了无穷维空间两个范数等价的\textbf{充分条件}:\par
\textbf{推论}$^\clubsuit $:设$\|\cdot\|$与$\|\cdot\|'$是同一个向量空间$X$上的两个具有以下性质的范数:空间$(X,\|\cdot\|)$与$(X,\|\cdot\|')$都是完备的,且存在常数$C$使得
$$
\|x\|'\leqslant C\|x\|,\forall x\in X
$$
则两个范数$\|\cdot\|$与$\|\cdot\|'$等价.\par
\textbf{Proof}:根据假设,恒等算子$T:(X,\|\cdot\|)\to (X,\|\cdot\|')$是双射且连续线性的.根据逆算子定理知,$T^{-1}:(X,\|\cdot\|')\to (X,\|\cdot\|)$存在且连续线性,从而有界.这表明存在常数$C'$使得
$$
\|x\|\leqslant C'\|x\|',\forall x\in X
$$
因此,两个范数是等价的.\par
\textbf{定义}$^\clubsuit $:给定两个集合$X$与$Y$,映射$T:D(T)\subset X\to Y$的图像$\operatorname{Gr}(T)$是如下定义的乘积$X\times Y$的子集
$$
\operatorname{Gr}(T):=\{(x,y)\in X\times Y:x\in D(T),y=Tx\}
$$
\qquad 若$X$与$Y$是赋范线性空间,一个映射$T:D(T)\subset X\to Y$被称作是闭的,是指其图像$\operatorname{Gr}(T)$在乘积空间$X\times Y$(装备积拓扑)中是闭的.\par
事实上,我们可以在乘积空间$X\times Y$上定义如下范数:对于$z=(x,y)\in X\times Y$,令
$$
\|z\|'=\|(x,y)\|'=\|x\|_{X}+\|y\|_{Y}
$$
由此,我们就使$X\times Y$成为了一个赋范线性空间,自然具备拓扑.\par
如果我们给定一个具体的算子,如何判定它是否是闭算子呢?利用定义的话,似乎不太方便.幸运的是,有下面的定理:\par
\textbf{定理}$^\clubsuit $:设$X,Y$是赋范线性空间,$T:D(T)\subset X\to Y$是线性算子.则$T$是闭线性算子的充分必要条件是:对任意的$\{x_n\}\subset D(T)$,满足$x_n\to x,Tx_n\to y$的时候,必有$x\in D(T)$且$Tx=y$.\par
\textbf{Proof}:若$T$是闭线性算子,则$\operatorname{Gr}(T)$是闭集,并对于任意$x_n\in D(T)$,当$x_n\to x,Tx_n\to y$时,有$(x_n,Tx_n)\to (x,y)$.因此$(x,y)\in \operatorname{Gr}(T)$,由$\operatorname{Gr}(T)$的定义可以知道,$x\in D(T),Tx=y$.\par
另一方面,若$(x_n,Tx_n)\in\operatorname{Gr}(T)$,且$(x_n,Tx_n)\to (x,y)$时一定有$x\in D(T),Tx=y$,从而$(x,y)=(x,Tx)\in\operatorname{Gr}(T)$.所以$\operatorname{Gr}(T)$是闭集,即$T$是闭线性算子.
\begin{tcolorbox}[colback=KuangNei,colframe=BianKuang,title=\textbf{Banach闭图像定理},sharpish corners]
    设$X$及$Y$是Banach空间,而$T\in\mathcal{L}(X;Y)$是闭线性算子.则$T\in \mathcal{L}(X;Y)$.
\end{tcolorbox}
\textbf{Proof}:在$X$上定义另一个范数
$$
\|x\|'=\|x\|_{X}+\|Tx\|_{Y},\forall x\in X
$$
根据范数$\|\cdot\|'$的定义,对于任何关于范数$\|\cdot\|'$的Cauchy列$\{x_n\}_{n=1}^{\infty}$,$\{x_n\}_{n=1}^{\infty}$和$\{Tx_n\}_{n=1}^{\infty}$分别是空间$X$及空间$Y$的Cauchy列.
由于$X,Y$都完备,故存在$x\in X$和$y\in Y$使得
$$
\lim\limits_{n\to\infty}x_n=x\text{在$X$中},\quad \lim\limits_{n\to\infty}Tx_n=y\text{在$Y$中}
$$
由假设,$T$是闭的,所以$y=Tx$.因此当$n\to\infty$时,
$$
\|x_n-x\|'=\|x_n-x\|_{X}+\|Tx_n-Tx\|_{Y}=\|x_n-x\|_{X}+\|Tx_n-y\|_{Y}\to 0
$$
这说明$(X,\|\cdot\|')$也是完备的.\par
易知$\|x\|\leqslant\|x\|'$对所有$x\in X$成立,这说明两个范数$\|\cdot\|$和$\|\cdot\|'$等价,故而存在常数$C$使得
$$
\|Tx\|_{Y}\leqslant \|x\|_{X}+\|Tx\|_{Y}=\|x\|'\leqslant C\|x\|,\forall x\in X
$$
所以线性算子$T$有界,故连续.\par
事实上,这个定理的逆定理也是正确的,因为任何连续映射的图像都是闭的.
\begin{tcolorbox}[colback=KuangNei,colframe=BianKuang,title=\textbf{实向量空间中的Hahn-Banach定理:解析形式},sharpish corners]
    设$X$是实向量空间,而$p$是$X$上的次线性泛函,即一个满足下述性质的函数$p:X\to \mathbb{R}$:
    $$
    \begin{aligned}
        p(t x) &=t p(x),\forall x\in X\text{且}\forall t>0\\
        p(x+y)&\leqslant p(x)+p(y),\forall x,y\in X
    \end{aligned}
    $$
    另设$Y$是$X$的子空间,$f:Y\to \mathbb{R}$是$Y$上的线性泛函,它满足
    $$
    f(x)\leqslant p(x),\text{对所有的}x\in Y
    $$
    则存在线性泛函$F:X\to \mathbb{R} $,它延拓$f$,即$\forall x\in Y,F(x)=f(x)$,且
    $$
    F(x)\leqslant p(x),\forall x\in X.
    $$
\end{tcolorbox}
\textbf{Proof}:定义集合
$$
P=\left\{h:D(h)\subset X\to\R\,\raisebox{-5.0ex}{\rule{0.9pt}{11.0ex}}\,
    \begin{aligned}
    &D(h)\text{是}X\text{的线性子空间},\\
    &h\text{是线性的},Y\subset D(h),\\
    &h\text{延拓}f,\text{且}h(x)\leqslant p(x),\text{对所有的}x\in D(h)
    \end{aligned}\right\}
$$
在$P$上我们定义序关系"$\prec $":
$$
h_1\prec h_2\Leftrightarrow D(h_1)\subset D(h_2),h_2(x)\text{延拓}h_1(x)
$$
因为$f\in P$,所以$P$非空.设$Q$为$P$的全序子集,记$Q=\{h_i\}_{i\in I}$.我们令
$$
D(h)=\bigcup_{i\in I}D(h_i),\quad h(x)=h_i(x)\text{若}x\in D(h_i)
$$
显然$h$是良定义的,并且$h\in P$.容易发现$h$是$Q$的一个上界,由Zorn引理可知,$P$有极大元,记为$F$.我们现在需要证明$D(F)=X$.\par
若$D(F)\neq X$,我们尝试能否将$F$再在$P$的意义下延拓.存在$x_0\in X\setminus D(F)$.定义$E=D(F)+\R x_0$.构造$E$上的线性泛函
$$
h(x+tx_0)=F(x)+t\alpha,\quad \forall x\in D(F),t\in \R
$$
事实上$\alpha=F(x_0)$.我们现在的问题是能否选取$\alpha$使得$h\in P$.即,要使$h\leqslant p$,只需要选取$\alpha$满足
$$
h(x+tx_0)=F(x)+t\alpha\leqslant p(x+tx_0),\quad \forall x\in D(F),t\in \R
$$
值得注意的是,对于固定的$x$,$p(x+tx_0)$是关于$t$的凸函数(不要忘了$p$是次线性泛函),而$F(x)+t\alpha$是关于$t$的线性函数.我们要让上面的式子成立,其实就是要让直线$F(x)+t\alpha$永远在凸函数$p(x+tx_0)$的下方.
又因为直线只需要两点确认,因此只需要验证在两个点处成立.这里我们取$t=\pm 1$,即
$$
\begin{cases}
    F(x)+\alpha \leqslant p(x+x_0),&\forall x\in D(F)\\
    F(x)-\alpha \leqslant p(x-x_0),&\forall x\in D(F)
\end{cases}
$$
这就是说,我们希望找到一个$\alpha$同时满足上面两个不等式,即
$$
\sup_{y\in D(F)}\{F(y)-p(y-x_0)\}\leqslant \alpha \leqslant \inf_{x\in D(F)}\{p(x+x_0)-F(x)\}
$$
这样的$\alpha$是存在的.因为的确有
$$
F(y)-p(y-x_0)\leqslant p(x+x_0)-F(x),\forall x,y\in D(F)
$$
上面的不等式成立需要借助$p$的次线性性质,由此有
$$
F(x)+F(y)\leqslant p(x+y)\leqslant p(x+x_0)+p(y-x_0)
$$
\qquad 通过上面的论证,我们发现了$h\in P$,且$F\prec h$.这个与$F$是$P$的极大元矛盾.因此$D(F)=X$.\par
我们知道对于对于一般数域上的赋范线性空间$X$来讲,$X'$上的范数可定义为
$$
\|f\|_{X'}=\sup_{x\in X,\|x\|=1}|f(x)|=\sup_{x\in X,\|x\|\leqslant 1}|f(x)|
$$
事实上,若$X$是实数域上的赋范线性空间,则$X'$上的范数可以定义为
$$
\|f\|_{X'}=\sup_{x\in X,\|x\|=1}|f(x)|=\sup_{x\in X,\|x\|\leqslant 1}|f(x)|=\sup_{x\in X,\|x\|\leqslant 1}f(x)
$$
你会发现最后把绝对值去掉了(对于复数域不可以这样).这是因为对于连续线性泛函$f\in X'$来讲,其范数就是$|f|$在单位球上的最大值.而单位球是对称的,若$\|x\|\leqslant 1$,那么必然有$\|-x\|\leqslant 1$.因此必然有
$$
\sup_{x\in X,\|x\|\leqslant 1}f(x)=\sup_{x\in X,\|x\|\leqslant 1}f(-x)
$$
又因为$f(x)=-f(-x)$,所以
$$
\inf_{x\in X,\|x\|\leqslant 1}f(x)=\inf_{x\in X,\|x\|\leqslant 1}-f(-x)=-\sup_{x\in X,\|x\|\leqslant 1}f(-x)=-\sup_{x\in X,\|x\|\leqslant 1}f(x)
$$
故
$$
\|f\|_{X'}=\sup_{x\in X,\|x\|=1}|f(x)|=\max\{\sup_{x\in X,\|x\|\leqslant 1}f(x),-\inf_{x\in X,\|x\|\leqslant 1}f(x)\}=\sup_{x\in X,\|x\|\leqslant 1}f(x)
$$
\qquad 现在我们可以给出几个推论:\par
\textbf{推论1}$^\clubsuit $:设$G\subset X$为线性子空间,若$g:G\to \R$为连续线性泛函,则存在$f\in X'$延拓$g$,且满足
$$
\|f\|_{X'}=\sup_{x\in G,\|x\|\leqslant 1}|g(x)|=\|g\|_{G'}
$$
\qquad \textbf{Proof}:$f$的存在性有Hahn-Banach定理保证,只需证明范数相等.显然
$$
\|f\|_{X'}=\sup_{x\in X,\|x\|\leqslant 1}|f(x)|\geqslant \sup_{x\in G,\|x\|\leqslant 1}|g(x)|=\|g\|_{G'}
$$
另一方面,在Hahn-Banach定理的证明中,我们可以选取$p(x)=\|g\|_{G'}\|x\|$.因此由$f(x)\leqslant p(x)$:
$$
\|f\|_{X'}=\sup_{x\in X,\|x\|\leqslant 1}f(x)\leqslant \|g\|_{G'}
$$
这就证明了结论.\par
\textbf{推论2}$^\clubsuit $:设$X$为实赋范线性空间,$\forall x_0\in X,x_0\neq 0$,存在$X$上的连续线性泛函$f$,使得$\|f\|_{X'}=1,f(x_0)=\|x_0\|$.\par
\textbf{Proof}:令$G=\R x_0$,显然$G$是$X$的线性子空间.在$G$上定义线性泛函
$$
f_0(\alpha x_0)=\alpha\|x_0\|
$$
事实上$f_0$还是连续的.对于$x\in G$,则$x=\alpha x_0$,故
$$
|f_0(x)|=|f_0(\alpha x_0)|=|\alpha|\|x_0\|=\|\alpha x_0\|=\|x\|
$$
这说明了$\|f_0\|_{X'}=1$,即$f$有界,所以连续.并且当$\alpha=1$时,$f_0(x_0)=\|x_0\|$.\par
有推论1可知存在$f\in X'$,使得$\|f\|_{X'}=1,f(x_0)=\|x_0\|$.\par
下面我们将给出Hahn-Banach定理的另一个重要的推论,它给出了赋范线性空间可分性的一个充分条件,为此我们需要先证明关于上面的推论2的推广(在下述定理中令$Y=\{0\}$就得到推论2),这个推广证明了:如果$Y$是赋范线性空间的一个闭的真子空间,那么一定存在一个$X$上的非零连续泛函,在$Y$上为零.\par
\textbf{定理}$^\clubsuit $:设$X$是赋范线性空间,$Y$是$X$的一个闭的真子空间.给定任意$x\in X-Y$,一定存在$f_x\in X'$使得
$$
f_x(y)=0,\quad \text{对所有}y\in Y,\quad f_x(x)=\inf\limits_{y\in Y}\|x-y\|>0\text{且}\|f_x\|=1
$$
\qquad \textbf{Proof}:定义$X$的子空间$Z$以及函数$f:Z\to\mathbb{K}$如下
$$
\begin{aligned}
    Z &= \mathbb{K}x+Y\\
    f(\alpha x+y)& =\alpha d,\forall \alpha \in\mathbb{K},d=\inf\limits_{y\in Y}\|x-y\|
\end{aligned}
$$
那么我们定义了一个$Z$上的线性泛函,且满足$f(x)=d,f(y)=0,\forall y\in Y$.断言,$f$连续且$\|f\|_{Z'}=1$.\par
由$d$的定义,对$\forall \alpha x+y\in Z,\alpha\neq 0$有
$$
\|\alpha x+y\|=|\alpha|\|x+\alpha^{-1}y\|\geqslant |\alpha||d|
$$
于是$|f(\alpha x+y)|=|\alpha||d|\leqslant \|\alpha x+y\|$.事实上,这个不等式对$\alpha=0$也是成立的.这就表明$\|f\|_{Z'}\leqslant 1$.\par
另一方面,依然由$d$的定义,对$\forall \varepsilon>0$,存在$y_{\varepsilon}\in Y$使得$\|x-y_{\varepsilon}\|<d+\varepsilon$.由于$x-y_{\varepsilon}\in Z$且$f(x-y_{\varepsilon})=f(x)=d$.所以
$$
\|f\|_{Z'}=\sup_{z\neq 0}\frac{|f(z)|}{\|z\|}\geqslant \frac{|f(x-y_{\varepsilon})|}{\|x-y_{\varepsilon}\|}\geqslant\frac{d}{d+\varepsilon}
$$
由$\varepsilon>0$的任意性可以得出$\|f\|_{Z'}\geqslant 1$.所以$\|f\|_{Z'}= 1$,即$f\in Z'$.\par
由Hahn-Banach延拓定理知,存在$f_x\in X'$使得
$$
f_x(z)=f(z),\forall z\in Z\text{且}\|f_x\|=\|f\|=1
$$
特别的有$f_x(x)=f(x)=d$且$f_x(y)=0,\forall y\in Y$.\par
容易发现,定理中的闭条件只是为了保证$d>0$,即使$Y$不是闭的,该定理对任意满足$\inf\limits_{y\in Y}\|x-y\|>0$的$x\in X-Y$依然成立.
\begin{tcolorbox}[colback=KuangNei,colframe=BianKuang,title=\textbf{赋范线性空间可分性的充分条件},sharpish corners]
如果赋范线性空间$X$的对偶空间$X'$是可分的,则$X$也是可分的.
\end{tcolorbox}
\textbf{Proof}:
























\newpage
设$X$为一赋范线性空间,所谓$X$上的超平面$H$,就是$X$的具有下面形式的子集$H$:
$$
H=\{x\in X:f(x)=\alpha\}
$$
其中$f:X\to\R$为非零线性泛函(注意,我们没有说它连续),$\alpha\in R$为给定常数.\par
设$H=[f=\alpha]$为$X$上的超平面,$A\subset X,B\subset X$.如果
$$
f(x)\leqslant \alpha,\quad \forall x\in A,\,\text{且}\,f(y)\geqslant \alpha,\quad \forall y\in B
$$
我们就称超平面$H=[f=\alpha]$分离$A$和$B$.如果还存在$\varepsilon>0$使得
$$
f(x)\leqslant \alpha-\varepsilon,\quad \forall x\in A,\,\text{且}\,f(y)\geqslant \alpha+\varepsilon,\quad \forall y\in B
$$
就称超平面$H=[f=\alpha]$严格分离$A$和$B$.\par
给定线性空间$X$的子集$A$,如果它满足
$$
tx+(1-t)y\in A,\quad \forall x,y\in A,t\in [0,1]
$$
则称$A$为凸集.
\begin{tcolorbox}[colback=KuangNei,colframe=BianKuang,title=\textbf{Hahn-Banach定理的第一几何形式:凸集的分离},sharpish corners]
设$A$与$B$为实赋范向量空间 $X$ 的两个非空子集,其满足以下性质:
$$
A\text{是开凸集};\quad B\text{是凸集};\quad A\cap B=\emptyset
$$
那么存在非零的 $f \in X'$ 及 $\alpha\in\mathbb{R}$使得
$$
f(x)\leqslant \alpha\leqslant f(y),\forall x\in A,y\in B
$$
\end{tcolorbox}
\textbf{Proof}:证明分成三个部分:\par
\textbf{Part1.构造一个次线性泛函(称为Minkowski泛函)}.\par
设$C$是$X$的包含原点的非空凸开子集.定义的函数$p:X\to[0,\in fty)$如下
$$
p(x)=\inf\left\{\beta>0:{\frac{x}{\beta}}\in C\right\}
$$
下面证明$p(x)$具有如下性质:\\
(a).存在一个常数$M$使得$0\leqslant p(x)\leqslant M\|x\|,\forall x\in X$;\\
(b).\,$C=\{x\in X;p(x)<1\}$;\\
(c).\,$p(x)$是一个次线性泛函.\par
因为$C$是开的且包含原点,故存在$r>0$使$B(0;r)\subset C$.给定任意$x\in X$及任意的$\beta>{\frac{\|x\|}{r}}$,则$\frac{x}{\beta}\in B(0;r)$故$\frac{x}{\beta}\in C$.所以
$$
p(x)=\inf\left\{\beta>0:{\frac{x}{\beta}}\in C\right\}\leqslant \inf\left\{\beta>{\frac{\|x\|}{r}}\right\}={\frac{\|x\|}{r}}.
$$
因此有$0\leqslant p(x)\leqslant M\|x\|$对所有$x\in X$其中$M:={\frac{1}{r}}$.\par
给定任意$x\in C$.因$C$是开的,存在$\delta>0$使$(1+\delta)x\in C$.这就有
$$
p(x)=\inf\left\{\beta>0:{\frac{x}{\beta}}\in C\right\}\leqslant {\frac{1}{1+\delta}}<1.
$$
反之,设$x\in X$且$p(x)<1$.这意味着存在$0<\beta<1$使得$\frac{x}{\beta}\in C$.因为$C$是凸的且包含原点,故$x=\left(\beta\frac{x}{\beta}+(1-\beta)0\right)\in C$.所以$C=\{x\in X;p(x)<1\}$.\par
给定任意$\alpha>0$及任意的$x\in X$
$$
\begin{aligned}
    p(\alpha x) &= \inf\left\{\widetilde{\beta}>0:\frac{\alpha x}{\widetilde{\beta}}\in C\right\}\\
                &= \inf\left\{\alpha\beta>0:\frac{\alpha x}{\alpha\beta}=\frac{x}{\beta}\in C\right\}=\alpha p(x)
\end{aligned}
$$
此外,$p(0)=0$所以$p(\alpha x)=\alpha p(x)$对所有$\alpha\geqslant0$及所有$x\in X$成立.\par
最后,设两点$x,y\in X$及$\varepsilon>0$是给定的.由上述性质有
$$
p\left(\frac{x}{p(x)+\varepsilon}\right)=\frac{p(x)}{p(x)+\varepsilon}<1.
$$
因此$\left(\frac{x}{p(x)+\varepsilon}\right)\in C$,同理$\left(\frac{y}{p(y)+\varepsilon}\right)\in C$这样，$c$的凸性意味着
$$
\left(\frac{\mu}{p(x)+\varepsilon}x+{\frac{1-\mu}{p(y)+\varepsilon}}y\right)\in C
$$
对所有$0<\mu<1$成立.特别地选取$\mu=\frac{p(x)+\varepsilon}{p(x)+p(y)+\varepsilon}$就有
$$
\frac{1}{p(x)+p(y)+2\varepsilon}(x+y)\in C.
$$
这就得到(这是因为$p\left(\frac{1}{p(x)+p(y)+2\varepsilon}(x+y)\right)<1$)
$$
p(x+y)<p(x)+p(y)+2\varepsilon.
$$
因为$\varepsilon>0$是任意的,所以$p(x+y)\leqslant p(x)+p(y)$.\par
这就证明了函数$p$具有所列出的所有性质.\par
\textbf{Part2.设$C$是$X$的一个非空凸开子集(这里没有说$C$是否包含原点),而$y_{0}\not\in C$.则存在$f\in X'$使得}
$$
f(x)<f(y_0),\forall x\in C
$$
先假定原点$0\in C$.设$p:X\to[0,\in fty)$为由Part1中所定义的函数,$Y=\mathrm{Span}\{y_{0}\}$(因$0\in $$C$,而$y_{0}\not\in C$,故$y_{0}$是非零向量),而$f_{0}:Y\to\mathbb{R}$为如下定义的线性泛函:对每个$\alpha\in \mathbb{R},f_{0}(\alpha y_{0})=\alpha$.当$\alpha\geqslant0$时,由Part1,$y_{0}\not\in C$,则$p(y_{0})\geqslant1$,故
$$
f_0(\alpha y_0)=\alpha\leqslant \alpha p(y_0)=p(\alpha y_0)
$$
当$\alpha\leqslant 0$时,由$p(y)\geqslant0$对所有$y\in X$,所以
$$
f_0(\alpha y_0)=\alpha\leqslant 0\leqslant p(\alpha y_0)
$$
上面的论证表明
$$
f_0(y)\leqslant p(y),\forall y\in Y
$$
因为$p$是一个次线性泛函,根据实向量空间中的Hahn-Banach定理,存在一个非零线性泛函$f:X\to\mathbb{R}$使得
$$
f(y_0)=f_0(y_0)=1,\text{且}f(x)\leqslant p(x),\forall x\in X
$$
又对所有$x\in X$成立不等式$p(x)\leqslant M\|x\|$,这意味着
$$
f(x)\leqslant p(x)\leqslant M\|x\|\quad\text{及}\quad -f(x)=f(-x)\leqslant p(-x)\leqslant M\|x\|
$$
所以线性泛函$f$是连续的,即$f\in X'$.此外,
$$
f(x)\leqslant p(x)<1=f_0(y_0)=f(y_0),\forall x\in C
$$
所以如果原点$0\in C$,上述断言就得证.\par
现在假定原点$0\notin C$.任意选取点$x_{0}\in C$,令
$$
\widetilde{C}=\{(x-x_0)\in X:x\in C\},\tilde{y}_0=y_0-x_0
$$
因为$0\in {\widetilde{C}}$且$\widetilde{y}_{0}\notin \widetilde{C}$，前面的推导说明,存在$f\in X'$使得
$$
f(\tilde{x})<f(\tilde{y}_0),\forall\tilde{x}\in \widetilde{C}
$$
由于$f$是线性的,因此$f(x)<f(y_{0})$对所有$x\in C$成立.\par
\textbf{Part3.完成证明}.\par
最后,设$A$是$X$的非空凸开子集.$B$是$X$的非空凸子集使得$A\cap B=\emptyset$.\par
定义集合
$$
C_1=\bigcup_{y\in B}\{(x-y)\in X:x\in A\},
$$
作为开集的并,它是开集.实际上它也是凸的:设$z_{i}=(x_{i}-y_{i})\in C_1,$其中$x_{i}\in A,y_{i}\in B,i=1,2$.因为两个集合$A$与$B$都是凸的,则
$$
(\mu x_{1}+(1-\mu)x_{2})\in A,(\mu y_{1}+(1-\mu)y_{2})\in B,\forall0<\mu<1
$$
因此
$$
\mu z_{1}+(1-\mu)z_{2}=\left((\mu x_{1}+(1-\mu)x_{2})-(\mu y_{1}+(1-\mu)y_{2})\right)\in C_1
$$
最后,因为$A\cap B=\emptyset$，所以$0\notin C_1$
根据Part2,存在非零的$f\in X'$使得$f(z)<f(0)=0$对所有$z\in C_1$成立,又$z=x-y$,其中$x\in A,y\in B$,即有
$$
f(x)<f(y),\text{对所有的}x\in A\text{及所有的}y\in B
$$
因此存在$\alpha\in \mathbb{R}$使得
$$
f(x)\leqslant \alpha\leqslant \inf\limits_{y\in B}f(y),\quad \forall x\in A
$$
\begin{tcolorbox}[colback=KuangNei,colframe=BianKuang,title=\textbf{Hahn-Banach定理的第二几何形式:凸集的严格分离},sharpish corners]
设$A$与$K$为实赋范向量空间$X$的两个非空凸子集,其满足以下性质:
$$
A\text{是闭集};\quad K\text{是紧集};\quad A\cap K=\emptyset
$$
那么存在非零的$f\in X',\gamma \in\mathbb{R}$及$\varepsilon>0$使得
$$
f(x)\leqslant \gamma-\varepsilon <\gamma +\varepsilon\leqslant f(y),\forall x\in A,y\in K
$$
\end{tcolorbox}
\textbf{Proof}:对任意 $r > 0$ ,设
$$
A(r)=\bigcup_{x\in A}B(x,r),\quad K(r)=\bigcup_{y\in K}B(y,r).
$$
则集合$A(r)$与$K(r)$都是非空,凸而且开的.此外,存在$r_{0}>0$,对所有$r\leqslant r_{0}$使得$A(r)\cap K(r)=\emptyset$.为说明这一结论,假定其不成立,于是对于任意$r>0$,存在$r'\leqslant r$使得$A(r')\cap K(r')\neq\emptyset$.取$r=\frac{1}{n}$,则存在$r_n\leqslant \frac{1}{n}$使得
$$
A(r_n)\cap K(r_n)\neq\emptyset
$$
设$z_n\in A(r_n)\cap K(r_n)$,根据$A(r_n),K(r_n)$的定义,则存在$x_n\in A(r_n)\subset A$及$y_n\in K(r_n)\subset K$满足
$$
\|z_n-x_n\|<r_n,\quad\|z_n-y_n\|<r_n
$$
令$v_n=z_n-x_n,w_n=z_n-y_n$,则$\|v_n\|<r_n,\|w_n\|<r_n$,并且
$$
x_n+v_n=z_n=y_n+w_n\\
$$
当$n\to \infty$时,$v_n\to 0,w_n\to 0$.\par
将上面的论述总结起来就是:存在$x_{n}\in A,v_{n}\in X,y_{n}\in K,w_{n}\in X,n\geqslant 1$,使得
$$
\begin{aligned}
&x_n+v_n=z_n=y_n+w_n,\quad n\geqslant 1\\
&v_n\to 0,w_n\to 0,\quad \text{当}n\to \infty\text{时}
\end{aligned}
$$
因集合$K$是紧的,则存在一个在$K$中收敛的子列$(y_{n_k})_{k=1}^{\infty}$.所以子列$(x_{n_k})_{k=1}^{\infty}$也在$X$中收敛,且由$A$是闭的,则$\lim_{k\to\infty}x_{n_k}\in A$．因此$\lim_{k\to\infty}x_{n_k}=\lim_{k\to\infty}y_{n_k}\in A\cap K$,这与假设矛盾.\par
事实上,对所有的$x\in X$,只要取$v\in X$满足$\|v\|<r_0$,则$x+v\in A(r_0)$.并且$A(r_0)$中的元素也一定能写成$x+v$的形式.这对$K(r_0)$也一样.\par
所以,根据Hahn-Banach定理的第一几何形式,两个集合$A(r_{0})$与$B(r_{0})$被一个超平面分,，即存在非零的$f\in X'$及$\gamma\in \mathbb{R}$使得
$$
f(x+v)\leqslant\gamma\leqslant f(y+w)
$$
对所有 $x\in A,y\in K$及所有$v,w\in X,\|v\|<r_0,\| w \|<r_{0}$成立.我们需要注意到下面的事实:$f$是连续的线性泛函,因此在闭球上一定可以达到最值,并且在开球上的上确界就等于在闭球上的上确界,这是下面利用范数放缩的关键.\par
固定$x\in A$:
$$
\gamma\geqslant \sup_{\|v\|<r_0}f(x+v) =f(x)+\sup_{\|v\|<r_0}f(v)=f(x)+\sup_{\|v\|\leqslant r_0}f(v)=f(x)+r_0\|f\|
$$
固定$y\in K$:
$$
\gamma\leqslant \inf_{\|w\|<r_0}f(y+w) =f(y)+\inf_{\|w\|<r_0}f(w)=f(y)+\inf_{\|w\|\leqslant r_0}f(w)=f(x)-r_0\|f\|
$$
令$\varepsilon=r_0\|f\|$,就能得到
$$
f(x)\leqslant \gamma-\varepsilon<\gamma+\varepsilon\leqslant f(y),\quad \forall x\in A,y\in K
$$
\newpage
\subsection{对偶空间与对偶算子}
\subsubsection{对偶空间}
设$X$是一赋范线性空间,我们知道$X$上所有连续线性泛函构成一线性空间$B(X,\mathbb{K})$,我们记为$X'$,称为$X$的对偶空间.如果在$X'$赋以算子范数
$$
\|f\|_{X'}=\sup_{x\in X,\|x\|=1}|f(x)|=\sup_{x\in X,\|x\|\leqslant 1}|f(x)|
$$
那么$X'$为Banach空间.对于$f\in X',x\in X$,一般我们将$f(x)$记为$\langle f,x\rangle$.如果继续定义下去,$(X')'=X''=B(X',\mathbb{K})$称为$X$的二次对偶空间.一个很自然的想法是$X$是否等于$X''$?现在给定(这时候固定了)$x\in X$,定义$X'$到$\mathbb{K}$上的连续线性泛函
$$
\begin{aligned}
    J_{x}:&X'\to \mathbb{K}\\
          &f\mapsto \langle f,x\rangle
\end{aligned}
$$
我们现在需要明白一点$J_{x}$是$X''$中的一个元素,并且它是依赖于$x$的,此时$J_{x}(f)=\langle f,x\rangle$中,变量是$f$,而不是$x$.按我们上面的记号也可以写成
$$
\langle J_{x},f\rangle =\langle f,x\rangle
$$
也就是说,任给一个$x$,都可以找到一个$X'$上的线性泛函$J_{x}$.这说明我们定义了一个从$X$到$X''$的映射,我们称之为\textbf{典则嵌入(自然映射)},用$J$来表示.即
$$
\begin{aligned}
    J:&X\to X''\\
      &x\mapsto J_{x}
\end{aligned}
$$
显然$J$是线性的,并且还是一个等距映射,即对所有的$x\in X,\|J_{x}\|_{X''}=\|x\|_{X}$,事实上
$$
\|J_{x}\|=\sup_{f\in X',\|f\|_{X'}=1}|\langle J_{x},f\rangle|=\sup_{f\in X',\|f\|_{X'}=1}|\langle f,x\rangle|=\|x\|_{X}
$$
上面等式中乍一看令人大意的地方在最后一个等号,它很像算子范数的定义.我们回忆以下算子范数的定义:$\forall f\in X'$
$$
\|f\|=\sup_{x\in X,\|x\|=1}|\langle f,x\rangle|
$$
再观察一下你就会发现不同,这似乎是关于$X$与$X'$范数的逆公式,它的合理性该如何解释呢?\par
\textbf{范数的逆公式\hypertarget{范数的逆公式}{}}$^\clubsuit $:设$X$是赋范线性空间,则对于$\forall x\in X$,成立
$$
\|x\|=\sup_{f\in X',\|f\|=1}|\langle f,x\rangle|
$$
\qquad \textbf{Proof}:从$f\in X'$,有$|\langle f,x\rangle|\leqslant \|f\|\cdot\|x\|$,所以
$$
\|x\|\geqslant \sup_{f\in X',\|f\|=1}|\langle f,x\rangle|
$$
对于另一个方向的不等式,$\forall x\neq 0,x\in X$,根据Hahn-Banach定理,存在$f\in X'$使得$\|f\|=1$且$\langle f,x\rangle=f(x)=\|x\|$.所以
$$
\sup_{f\in X',\|f\|=1}|\langle f,x\rangle|\geqslant|\langle f,x\rangle|=\|x\|
$$
从两个反向不等式就得到了结论,当然$x=0$的情况十分平凡,不值得赘述.\par
由$J$是等距映射,很快就可以得到$J$还是单射,但未必是满射.这说明$J$实际上是$X$与$X''$的一个子空间$J(X)$线性同构(还是保范的).那么在这种意义下,$X$可以看作是$X''$的一个子空间.\par
特别的,如果$J$还是满射,那么$X$和$X''$同构,也就是此时我们可以将$X$与$X''$看作用一个空间,这恰恰符合我们的直觉.并且这个时候我们称$X$为自反空间.\par
值得注意的是,作为一个对偶空间,自反空间肯定是完备的.\par
下面的定理给出了几个自反空间的例子:\par
\textbf{自反空间的例子}$^\clubsuit $:下述Banach空间是自反的:\par
(a).有限维赋范线性空间;\par
(b).自反Banach空间的任一闭子空间;\par
(c).自反Banach空间的对偶空间.\par
\textbf{Proof}:我们需要一个一个证明.\par
\textbf{(a)的证明}.设$X$是有限维赋范线性空间.给定$X$中的一组基$\{e_i\}^n_{i=1}$,利用$e'_j(e_i)=\delta_{ij},1\leqslant i,j\leqslant n$,定义了$X'$中的一组基$\{e_j'\}^n_{j=1}$.同理,利用$e''_k(e'_j)=\delta_{j,k},1\leqslant j,k\leqslant n$,定义了$X''$中的一组基.定义典则嵌入$J:X\to X''$如下
$$
J\left(\sum_{i=1}^{n}x_ie_i\right)=\sum^n_{i=1}e''_i,\quad \forall x=\sum_{i=1}^{n}x_ie_i\in X
$$
则$J$是满射.\par
\textbf{(b)的证明}.设$Y$是自反Banach空间$X$的闭子空间.我们只需要证明典则嵌入$J_{Y}:Y\to Y''$是满射即可,即我们要证明,给定任意的$y''\in Y''$,存在$y\in Y$使得
$$
y''(y')=y'(y),\quad \forall y'\in Y'
$$
\qquad 设$y''\in Y''$是给定的.则线性泛函
$$
\begin{aligned}
    x'':X' &\to \mathbb{K}\\
        x' &\to x''(x')=y''(x'|_{Y})
\end{aligned}
$$
是连续的.这是因为
$$
|x''(x')|=|y''(x'|_{Y})|\leqslant \|y''\|\|x'|_{Y}\|\leqslant \|y''\|\|x\|,\quad \forall x'\in X'
$$
所以$x''\in X''$.由于$X$是自反的.故存在$y\in X$使得
$$
x'(y)=x''(x')=y''(x'|_{Y}),\quad \forall x'\in X'
$$
根据上式,也可以看出,对于所有使得$x'|_{Y}=0$的那些$x'\in X'$,一定有$x'(y)=0$.已知$Y$是闭的,如果$y\notin Y$,一定存在$x'\in X'$使得$x'|_{Y}=0$,但$x'(y)\neq 0$.这就说明了$y\in Y$.\par
给定任意的$y'\in Y'$,$x'\in X'$是$y'$的任一延拓(由Hahn-Banach延拓定理保证这样的延拓存在),则$y'=x'|_{Y}$
$$
y''(y')=y''(x'|_{Y})=x'(y)=y'(y)
$$
\qquad \textbf{(c)的证明}.给定一个自反的Banach空间,设$J:X\to X''$为$X$到$X''$的典则嵌入,而$J':X'\to (X')''$为$X'$到$(X')''$的典则嵌入.我们要证明$J'$是满射,即对于任意$x'''\in (X')''$,都存在$x'\in X'$,使得$J'_{x'}=x'''$.\par
给定任意$x'''\in (X')''$,定义映射
$$
\begin{aligned}
    x':X &\to \mathbb{K}\\
        x &\to x'(x)=x'''(J_{x})
\end{aligned}
$$
注意$J_x\in X''$,这个定义是有意义的.又因为
$$
|x'(x)|\leqslant \|x'''\|_{(X')''}\|J_x\|_{X''}=\|x'''\|_{(X')''}\|x\|_{X},\quad \text{对所有}x\in X
$$
故一定有$x'\in X'$.(接下来的事情就是在说明这个$x'$就是我们寻找的).此外$(X')''=(X'')'$.由$x',J$的定义及$J$的双射性质得到
$$
x'''(x'')=x''(J_x)=x'(x)=J_x(x')=x''(x'),\quad x''=J_x\in X''
$$
再根据$J'$的定义,有
$$
x'''=J'_{x'}
$$
这说明了$J'$是满射.\par
\textbf{注}:上面的证明中出现了一个等式
$$
x'''(x'')=x''(J_x)=x'(x)=J_x(x')=x''(x'),\quad x''=J_x\in X''
$$
第一个等号是利用$J$的满射,对于任意$x''\in X''$,都存在$x\in X$使得$x''=J_x$;第二个等号就是用$x'$的定义了;第三个等号则是利用$J_x$的定义;最后一个等号和第一个等号一样.\par
上面只给了三个自反空间的例子,可以发现在证明中,我们是严格按照自反空间的定义--典则嵌入是满射--来证明的.定义当然不是证明的唯一办法,毕竟有时候用定义证还是很复杂的.下面介绍的Milman-Pettis定理说明,当我们证明了一个Banach空间一致凸,那它一定是自反的.我们将用这个定理证明$l^{p}$和$L^{p}(\Omega),1<p<\infty$是自反的.\par
为此,我们需要引入一致凸的概念:给定赋范线性空间$(X,\|\cdot\|)$,如果$\forall \varepsilon >0$,存在$\delta=\delta(\varepsilon)>0$,使得对一切满足
$$
\|x\|=\|y\|=1,\|x-y\|\geqslant \varepsilon
$$
的向量$x,y\in X$,都有
$$
\left\|\frac{x+y}{2}\right\|\leqslant 1-\delta
$$
则称$X$是一致凸的.事实上,在Wikipedia上,你会看到它会将$\varepsilon$限制为$0<\varepsilon\leqslant 2$.这是因为在单位球面上两点的距离不会超过2.另外还可以将单位球面的要求($\|x\|=\|y\|=1$)减弱为单位球内($\|x\|=\|y\|\leqslant 1$).\par
\begin{tcolorbox}[colback=KuangNei,colframe=BianKuang,title=\textbf{Milman-Pettis定理},sharpish corners]
一致凸的Banach空间一定是自反的.
\end{tcolorbox}
在有限维赋范线性空间我们知道,任一有界序列必然包含一个收敛子序列.然而这一论断对无限维空间未必成立.不过当我们引入弱收敛的概念时,将证明:\textbf{自反Banach空间中任一有界序列一定包含一个弱收敛的子序列}.而收敛性是和空间上定义的拓扑相关的,上面说明我们可以改变空间的拓扑来改善收敛.\par
现在来考虑另一件事情:拓扑和连续的关系.设有两个拓扑空间$(X,\mathcal{T}_X)$和$(Y,\mathcal{T}_Y)$,映射$f:X\to Y$称为连续的,当且仅当
$$
\text{对于任意的开集}V\subset Y,f^{-1}\in \mathcal{T}_X
$$
我们现在来改变$X$或者$Y$上的拓扑,观察连续性的变化情况.\par
1.如果$\mathcal{T}_X$变强(也就是它的开集更多),那么要求$f^{-1}(V)$是开集就更容易满足;反之,如果$\mathcal{T}_X$变弱(也就是它的开集更少),那么要求$f^{-1}(V)$是开集就更难满足(因为开集太少了,可供选择的开集不多);\par
2.如果$\mathcal{T}_Y$变强(也就是它的开集更多),那么要求$f^{-1}(V)$是开集就难满足,因为此时我们要验证更多$V$的原像;反之,如果$\mathcal{T}_X$变弱(也就是它的开集更少),那么要求$f^{-1}(V)$是开集就更容易满足,因为这时候验证的$V$数量就少了.\par
我们早就知道连续性和拓扑也相关,1.和2.都通过改变拓扑,都使得连续性改变,这说明我们可以通过改变空间上的拓扑来改善空间上的函数的连续性.\par
既然如此,我们何不寻找$X$最强的拓扑,$Y$最弱的拓扑一劳永逸呢?这太没意思了,毕竟此时$\mathcal{T}_Y$中只有两个元素.另外,上面只涉及一个映射.如果我们此时设$X$就是一个普通的集合,以及一族到拓扑空间$Y_i$的映射族$\{\varphi_i\}_{i\in I}$.是否存在$X$上的一个拓扑$\mathcal{T}_X$满足:\\
(a).所有的映射$\varphi _i:X\to Y_i,i\in I$是连续的.\\
(b).$X$的这个拓扑是使(a)成立的最弱的拓扑,也就是对于任意$X$上使(a)成立的拓扑$\mathcal{T}'_X$,都有$\mathcal{T}_X\subset \mathcal{T}'_X$.\par
答案是的确存在,而且还唯一,这种拓扑称为initial topology.具体构造方式是以
$$
\left\{\bigcap_{j\in J}\varphi^{-1}_{j}(U_j):U_j\in\mathcal{T}_j,J\subset I\text{为有限集}\right\}
$$
为基.这个拓扑我们并不打算深入下去,我们承认就行.现在回到泛函分析中:设$X$是Banach空间,且$f\in X'$.当$f$取遍$X'$时,$\{\langle f,x\rangle\}_{f\in X'}$就是一族从$X$到$\R$的映射,自然也可以在$X$上诱导initial topology.于是我们引入下面的概念:\par
\textbf{弱拓扑}$^\clubsuit $:$X$上的弱拓扑$\sigma (X,X')$是$X$上使得所有$\langle f,x\rangle,f\in X'$都连续的最弱的拓扑.\par
\textbf{注}:你可能会很疑惑,既然$f\in X'$,当然有$\langle f,x\rangle$是连续的了,上面的定义不是多此一举么.事实上,我们刚开始是在X用范数诱导了一个拓扑,此时X'中的元素都是在这范数拓扑下连续的线性泛函,现在我们要寻找一个比范数拓扑弱的拓扑(更经济实惠),但不破坏原来的X'中线性泛函的连续性(不减少旧的连续线性泛函),
并且只保证原来X'已有的线性泛函的连续性(不添加新的连续线性泛函),我们要找到这个拓扑就是上面定义的弱拓扑.\par
下面,我们先引入弱收敛的概念,再说明,弱收敛其实就是弱拓扑上的收敛.\par
设$X$是赋范线性空间,$X'$表示其对偶空间.序列$\{x_n\}_{n=1}^{\infty}\subset X$称为在$X$中是弱收敛的,是指存在$x\in X$使得对每个$f\in X'$
$$
f(x_n)\to f(x),\quad \text{当}n\to\infty\text{时}
$$
这个$x$则称为序列$\{x_n\}_{n=1}^{\infty}$的弱极限.我们用$x_n\rightharpoonup x$来表示弱收敛.\par
与弱收敛相对的自然就是强收敛,事实上,强收敛就是按范数拓扑收敛:即当$n\to\infty$时,$\|x_n-x\|\to 0$.强收敛依然用$\to$来表示.\par
\textbf{命题}$^\clubsuit $:在赋范线性空间$X$中,弱拓扑$\sigma (X,X')$的一个拓扑基由如下集合组成
$$
V(x;f_1,\cdots,f_n;\varepsilon)=\bigcap_{i=1}^n\{y\in X:|\langle f_i,y-x\rangle|<\varepsilon\}
$$
其中$x\in X,f_1,\cdots,f_n\in X',\varepsilon>0$.\par
\textbf{Proof}:令$\mathcal{B}=\{V(x;f_1,\cdots,f_n;\varepsilon)\}$,我们需要证明的是$\mathcal{B}$构成弱拓扑$\sigma (X,X')$的拓扑基.在前面我们已经知道,弱拓扑$\sigma (X,X')$的一个拓扑基可以写成
$$
\mathcal{B}_0=\left\{\bigcap_{i=1}^{n}f^{-1}_i(U_i):f_i\in X',U_i\subset \mathbb{K}\text{为开集}\right\}
$$
也就是说我们已经有了一个拓扑基$\mathcal{B}_0$,我们证明$\mathcal{B}$其实也是一个拓扑基.利用拓扑基判定定理:对于任意$x\in X$和任意开集$U\in \sigma (X,X')$,若$x\in U$,则存在$B\in\mathcal{B}$,使得$x\in B\subset U$.\par
任意$x\in X$和任意开集$U\in \sigma (X,X')$,且$x\in U$,由于$\mathcal{B}_0$是$\sigma (X,X')$的拓扑基,根据拓扑基的定义,一定存在某个基元素
$$
B_0=\bigcap_{i=1}^{n}f^{-1}_i(U_i)\in B_0\text{且}x\in B_0\subset U
$$
其中每个$f_i\in X',U_i\subset \mathbb{K}$是开集且$f_i(x)\in U_i$.\par
由于$U_i$是开集,故存在$\varepsilon _i>0$使得
$$
(f_i(x)-\varepsilon_i,f_i(x)+\varepsilon_i)\subset U_i
$$
现在定义$\varepsilon=\min_{1\leqslant i\leqslant n}\varepsilon _i$.记
$$
V=V(x;f_1,\cdots,f_n;\varepsilon)=\bigcap_{i=1}^n\{y\in X:|\langle f_i,y-x\rangle|<\varepsilon\}
$$
则对任意$y\in V$,我们有
$$
|\langle f_i,y-x\rangle|<\varepsilon\leqslant\varepsilon_i
$$
这就说明$f_i(y)\in (f_i(x)-\varepsilon_i,f_i(x)+\varepsilon_i)\subset U_i$,即$y\in f^{-1}_{i}(U_i)$,所以$y\in B_0$.从而$V\subset B_0\subset U$.\par
根据$V$的定义就可以知道$x\in V$且$V\in \mathcal{B}$,所以$x\in V\subset U$.这就说明了$\mathcal{B}$也是拓扑基.\par
现在我们就可以证明下面的命题:\par
\textbf{命题}$^\clubsuit $:设$X$是赋范线性空间,序列$\{x_n\}$在$X$中弱收敛于$x\in X$,当且仅当$x_n\to x$在$X$的弱拓扑$\sigma (X,X')$下收敛.\par
\textbf{Proof}:证明分成两个部分:\par
\textbf{Part1.充分性.}\par
假设$x_n$在弱拓扑$\sigma (X,X')$下收敛到$x\in X$.根据定义,对任意弱拓扑下的开邻域$U$包含$x$,存在$N>0$使得当所有的$n>N$,满足$x_n\in U$.由上一个命题可知弱拓扑的基中元素都可以写成
$$
U=\bigcap_{i=1}^m\{y\in X:|\langle f_i,y-x\rangle|<\varepsilon\}
$$
因此,从某个$N$起,对所有的$n>N$和$i=1,\cdots,m$,有
$$
|\langle f_i,x_n-x\rangle|<\varepsilon
$$
我们需要知道$U$是依赖于$f_i$和$\varepsilon$的,因此我们任意选取$U$,其实就是任意选取$f_i$和$\varepsilon$.因为$f_i$和$\varepsilon$是任意选的,此时我们可以给定任意的$f\in X'$,任意$\varepsilon>0$,取$m=1$且$f_1=f$(这时我们就构造了一个特殊的邻域$U$,这个就像数分中数列极限一样,我们在题目中已知极限存在,通常都会利用$\varepsilon$的任意性选择特殊的$\varepsilon$),于是当$n>N$,
$$
|\langle f,x_n-x\rangle|<\varepsilon
$$
这说明
$$
\lim_{n\to \infty}\langle f,x_n\rangle=\langle f,x\rangle
$$
\qquad \textbf{Part2.必要性.}\par
假设
$$
\lim_{n\to \infty}\langle f,x_n\rangle=\langle f,x\rangle,\forall f\in X'
$$
要证明在弱拓扑$\sigma (X,X')$下$x_n\to x$,就是要证明:对任意弱拓扑中包含$x$的开邻域$U$,存在$N>0$,使得对所有$n>N$时,$x_n\in U$.\par
根据拓扑基的结构,任意包含$x$的开邻域$U$都可以写成
$$
U=\bigcap_{i=1}^m\{y\in X:|\langle f_i,y-x\rangle|<\varepsilon\}
$$
因为对每个$i$,都有
$$
\lim_{n\to \infty}\langle f_i,x_n\rangle=\langle f_i,x\rangle
$$
所以对每个$i$,对任意$\varepsilon>0$,都存在$N_i(\varepsilon)>0$,当$n>N_i(\varepsilon)$时,有
$$
|\langle f_i,x_n-x\rangle|<\varepsilon
$$
令$N=\max_{1\leqslant i\leqslant m}N_i(\varepsilon)$,当$n>N$时,对所有的$i=1,\cdots,m$一致满足
$$
|\langle f_i,x_n-x\rangle|<\varepsilon
$$
这就说明了当$n>N$时,$x_n\in U=\bigcap_{i=1}^m\{y\in X:|\langle f_i,y-x\rangle|<\varepsilon\}$.\par
我们知道,对于任意给定的$x\in X$,映射$J_x$是$X'$上的连续线性泛函.我们在$X'$也引入一种拓扑:弱*拓扑.\par
\textbf{弱*拓扑}$^\clubsuit $:$X'$上的弱*拓扑$\sigma (X',X)$是$X'$上使得所有映射$J_x:X'\to\mathbb{K}$都连续的最弱的拓扑.\par
类比弱拓扑,弱*拓扑$\sigma (X',X)$上的收敛就是所谓的弱*收敛:设$X$是赋范线性空间,$X'$是其对偶空间,序列$\{f_n\}\subset X'$称为在$X'$中是弱*收敛,是指如果存在$f\in X'$使得对每个$x\in X$
$$
f_n(x)\to f(x),\quad \text{当}n\to\infty\text{时}
$$
称$f$为序列$\{f_n\}$的弱*极限,我们用$f_n\xrightharpoonup{*} f$表示弱*收敛.\par
在有限维赋范线性空间我们知道,单位闭球是紧的,这个结论放在无限维就不一定成立了.不过在我们引入弱*拓扑之后将会得到下面的结论
\begin{tcolorbox}[colback=KuangNei,colframe=BianKuang,title=\textbf{Banach-Alaoglu定理 },sharpish corners]
$X'$中的单位闭球
$$
B_{X'}=\{f\in X':\|f\|\leqslant 1\}
$$
对于$X'$的弱*拓扑$\sigma(X',X)$是紧的.
\end{tcolorbox}
作为结束,下面的定理回答了开头的疑问.并且这个定理在线性泛函分析中很重要:第一条在变分法中确立强制且序列弱下半连续的泛函极小元的存在性起着重要作用,第二条给出了证明一个空间是自反的一种方法.
\begin{tcolorbox}[colback=KuangNei,colframe=BianKuang,title=\textbf{Banach-Eberlein-\v{S}mulian定理 },sharpish corners]
(a).在自反Banach空间中,任何有界序列一定包含一个弱收敛的子列.\\
(b).反之,如果一个Banach空间中每一个有界序列均包含一个弱收敛的子列,则该空间一定是自反的.
\end{tcolorbox}
这个定理的证明不太容易,就略了.

\subsubsection{对偶算子}
不妨设$X,Y$是赋范线性空间空间,$T\in B(X,Y)$.任给$x\in X,f\in Y'$.固定$f$和$T$,那么$f(Tx)=\langle f,Tx\rangle$是定义在$X$上的一个泛函(这里其实和$x$并无关系,就好像我们写函数通常写为$f(x)$一样,仅表明作用的对象).
如果固定$T$,那么对于$\forall f\in Y'$,根据上面,我们定义了一个依赖于$T$的作用在$X$上的线性映射$T^*$:
$$
\begin{aligned}
    T^*:&Y'\to X'\\
          &f\mapsto T^*(f)(x)=f(Tx)
\end{aligned}
$$
也就是说$T^*f(x)=f(Tx)\Leftrightarrow \langle T^*f,x\rangle =\langle f,Tx\rangle$.我们称$T^*$是$T$的对偶算子.下面的定理保证了我们上面的论述是正确的.
\begin{tcolorbox}[colback=KuangNei,colframe=BianKuang,title=\textbf{对偶算子},sharpish corners]
设$X$与$Y$是同一数域$\mathbb{K}$上的两个赋范线性空间.给定任意算子$T\in\mathcal{L}(X;Y)$,存在唯一的算子$T^*\in\mathcal{L}(Y';X')$(称为$T$的对偶算子),使得
$$
T^*f(x)=f(Tx),\text{对所有的}x\in X\text{及所有的}f\in Y'
$$
此外$\|T^*\|_{\mathcal{L}(Y';X')}=\|T\|_{\mathcal{L}(X;Y)}$
\end{tcolorbox}
\textbf{Proof}:给定$f\in Y'$,映射
$$
\begin{aligned}
    T^*f:&X\to \mathbb{K}\\
          &x\mapsto T^*f(x)=f(Tx)
\end{aligned}
$$
作为连续线性映射的复合,是一个连续线性泛函.\par
由这种方式定义的映射$T^*:Y'\to X'$是线性的.事实上,对任意的$f_1.f_2\in Y'$,
$$
\begin{aligned}
    T^*(f_1+f_2)(x) &= (f_1+f_2)(Tx)=f_1(Tx)+f_2(Tx)\\
                    &= T^*f_1(x)+T^*f_2(x)
\end{aligned}
$$
对所有的$x\in X$成立;另外,对任意的$\alpha\in\mathbb{K}$及$f\in Y'$,有
$$
\begin{aligned}
    T^*(\alpha f)(x) &= (\alpha f)(Tx)=\alpha (f(Tx))\\
                     &= \alpha (T^*f(x))=(\alpha T^*f)(x)
\end{aligned}
$$
对所有的$x\in X$成立.\par
再证映射$T^*:Y'\to X'$是连续的.这是因为
$$
\begin{aligned}
    \|T^*f\| &=\sup_{x\neq 0}\frac{|T^*f(x)|}{\|x\|} =\sup_{x\neq 0}\frac{|f(Tx)|}{\|x\|}\\
             &\leqslant \|T\|\|f\|,\quad \forall f\in Y'
\end{aligned}
$$
所以有$\|T^*\|\leqslant \|T\|$(这已经表明$T^*$是有界的了,自然连续).另一方面,根据之前证过的\hyperlink{范数的逆公式}{范数的逆公式},
$$
\begin{aligned}
    \|Tx\| &=\sup_{f\neq 0}\frac{|f(Tx)|}{\|f\|}=\sup_{f\neq 0}\frac{|T^*f(x)|}{\|f\|}\\
           &\leqslant\left(\sup_{f\neq 0}\frac{\|T^*f\|}{\|f\|}\right)\|x\|=\|T^*\|\|x\|
\end{aligned}
$$
对所有的$x\in X$成立.因此$\|T\|\leqslant\|T^*\|$.所以$\|T^*\|=\|T\|$.\par


 \subsection{泛函分析复习选题}
\textbf{Problem 1}$^\clubsuit $:设$(X,d)$是一度量空间,$A\subset X$非空.令
$$
f(x)=\inf\{d(x,y):y\in A\},\quad x\in X
$$
证明:$f$是$X$上的连续函数.\par
\textbf{Proof}:对于任意$x,x_0\in X,y\in A$,跟据三角不等式得到
$$
d(x,y)\leqslant d(x,x_0)+d(x_0,y)
$$
对$y$取下确界,得到
$$
f(x)\leqslant d(x,x_0)+f(x_0)
$$
交换$x$和$x_0$的位置,得到$f(x_0)\leqslant d(x,x_0)+f(x)$.这表明
$$
|f(x)-f(x_0)|\leqslant d(x,x_0)
$$
只要取$\delta$满足$0<\delta<\varepsilon$就能得到$f$的连续性.\par
\textbf{Problem 2}$^\clubsuit $:设$A$为度量空间$(X,d)$的闭子集.证明:必有一列开集$\{O_n\}$使得$A\subset O_n,n=1,\cdots,$且$\bigcap_{n=1}^{\infty}O_n=A$.\par
\textbf{Proof}:对于每个$n=1,2,\cdots$,定义
$$
O_n=\{x\in X:d(x,A)<1/n\}
$$
显然$\{O_n\}$是开集列,且$A\subset O_n$.下面证明$\bigcap_{n=1}^{\infty}O_n\subset A$.\par
对于任意$x\in \bigcap_{n=1}^{\infty}O_n$,那么对于每个$n$,都有$d(x,A)<\frac{1}{n}$.令$n\to\infty$得到$d(x,A)=0$.而$A$是闭集,因此$x\in A$.\par
\textbf{Problem 3}$^\clubsuit $:设$(X,d)$为度量空间,$F_1,F_2$为$X$中不相交的闭集.证明:\par
1.存在连续函数$f:X\to [0,1]$使得$f(F_1)=\{0\},f(F_2)=\{1\}$;\par
2.存在$X$中的两个不相交的开集$G_1,G_2$使得$G_1\supset F_1,G_2\supset F_2$.\par
\textbf{Proof}:先证明第一个.事实上令
$$
f(x)=\frac{d(x,F_1)}{d(x,F_1)+d(x,F_2)}
$$
则$f$是$X$上的连续函数.并且
$$
f(x)=
\begin{cases}
0, & x\in F_1\\[4pt]
1, & x\in F_2
\end{cases}
$$
记$G_1=\{x\in X:f(x)<\frac{1}{2}\},G_2=\{x\in X:f(x)>\frac{1}{2}\}$,显然$G_1,G_2$是两个不相交的开集,且$G_1\supset F_1,G_2\supset F_2$.\par
\textbf{Problem 4}$^\clubsuit $:设$(X,d)$为紧度量空间,$T$是$X$到自身的映射且满足
$$
d(Tx,Ty)< d(x,y),\forall x,y\in X
$$
证明:$T$有唯一的不动点.\par
\textbf{Proof}:令$\phi(x)=d(Tx,x),x\in X$.先证明这是一个连续函数.事实上,$\forall x_1,x_2\in X$,有三角不等式
$$
d(Tx_1,x_1)\leqslant d(Tx_1,Tx_2)+d(Tx_2,x_2)+d(x_2,x_1)\leqslant d(Tx_2,x_2)+2d(x_2,x_1)
$$
交换$x_1,x_2$的位置,得到
$$
|d(Tx_1,x_1)-d(Tx_2,x_2)|\leqslant 2d(x_1,x_2)
$$
这就说明$\phi$是紧度量空间$X$上的连续函数,因此可以取到最小值.不妨设$x_0\in X$使得$\phi(x_0)=\min_{x\in X}\phi(x)$.如果$\phi(x_0)>0$,那么
$$
\phi(x_0)=d(Tx_0,x_0)>d(T^2x_0,Tx_0)=\phi(Tx_0)
$$
这和$\phi(x_0)$是最小值矛盾,因此只能有$\phi(x_0)=0$.这就说明$Tx_0=x_0$,即$T$有不动点.至于唯一性,显然$T$是单射.\par
\textbf{Problem 5}$^\clubsuit $:设$T$是从完备度量空间$X$到自身的映射,$\theta\in (0,1)$是一个常数.若在开球$B(x_0,r)(r>0)$内满足$d(x_0,Tx_0)\leqslant r(1-\theta)$且
$$
d(Tx,Ty)\leqslant \theta d(x,y),\forall x,y\in B(x_0,r)
$$
又在闭球$S(x_0,r)=\{x\in X:d(x,x_0)\leqslant r\}$上,$T$连续.证明:$T$在$S(x_0,r)$中有不动点.\par
\textbf{Proof}:我们只需要证明$T$将闭球映射到自身即可.事实上$\forall y\in S(x_0,r)$,有
$$
\begin{aligned}
    d(Ty,x_0) &\leqslant d(Ty,Tx_0)+d(Tx_0,x_0)\\
              &\leqslant \theta d(y,x_0)+d(Tx_0,x_0)\\
              &\leqslant \theta r+r(1-\theta)=r
\end{aligned}
$$
注意到,完备度量空间的闭子空间是完备的,根据Banach不动点定理就能得到结论.\par
求算子范数你经常能遇到$\sgn$函数.这里都是利用:若$x\in\mathbb{R}$,那么$x\cdot\sgn x=|x|$;若$x\in\mathbb{C}$,那么$x\cdot\overline{\sgn x}=|x|$.\par
\textbf{Problem 6}$^\clubsuit $:设无穷矩阵$(a_{ij}),i,j=1,2,\cdots$满足$\sup_{i\geqslant 1}\sum_{j=1}^{\infty}|a_{ij}|<\infty$.作$l^{\infty}$到$l^{\infty}$中的算子如下:若$x=(\xi_1,\xi_2,\cdots),y=(\eta_1,\eta_2,\cdots),Tx=y$,则
$$
\eta_i=\sum_{j=1}^{\infty}a_{ij}\xi_j,i=1,2,\cdots
$$
证明:$T$是$l^{\infty}$到$l^{\infty}$中的有界线性算子,并求$\|T\|$.\par
\textbf{Proof}:首先证明$T$是线性的.对于任意$x_1=(\xi_1^{(1)},\xi_2^{(1)},\cdots),x_2=(\xi_1^{(2)},\xi_2^{(2)},\cdots)\in l^{\infty}$,以及任意$\alpha,\beta\in \R$,则$T(\alpha x_1+\beta x_2)$的第$i$个分量为
$$
\eta_i=\sum_{j=1}^{\infty}a_{ij}(\alpha \xi_j^{(1)}+\beta \xi_j^{(2)})=\alpha \sum_{j=1}^{\infty}a_{ij}\xi_j^{(1)}+\beta \sum_{j=1}^{\infty}a_{ij}\xi_j^{(2)}
$$
其中关于级数的线性性质是因为我们是在$l^{\infty}$中考虑的(每个分量都是级数).下面我们求$\|T\|$.\par
事实上
$$
\begin{aligned}
    \|Tx\|_{l^{\infty}} &= \sup_{i\geqslant 1}|\eta_i|=\sup_{i\geqslant 1}\left|\sum_{j=1}^{\infty}a_{ij}\xi_j\right|\\
                       &\leqslant \sup_{i\geqslant 1}\sum_{j=1}^{\infty}|a_{ij}||\xi_j|\\
                       &\leqslant \left(\sup_{i\geqslant 1}\sum_{j=1}^{\infty}|a_{ij}|\right)\|x\|_{l^{\infty}}
\end{aligned}
$$
这就得到$\|T\|\leqslant \sup_{i\geqslant 1}\sum_{j=1}^{\infty}|a_{ij}|$.对每个固定的$i$,取$x^{(i)}=(\sgn a_{i1},\sgn a_{i2},\cdots)$,则$\|x^{(i)}\|_{l^{\infty}}=1$.并且
$$
\|T\|=\sup_{\|x\|=1}\|Tx\|_{l^{\infty}}\geqslant \|Tx^{(i)}\|_{l^{\infty}}\geqslant \sum_{j=1}^{\infty}|a_{ij}|
$$
再对上式取上确界就可以得到$\|T\|\geqslant \sup_{i}\sum_{j=1}^{\infty}|a_{ij}|$.这就证明了$T$是有界线性算子,并且求出了$T$的范数.\par

\textbf{Problem 7}$^\clubsuit $:设$X$是赋范线性空间,$V$是$X$中的完备子空间.若商空间$X/V$是完备的.证明$X$也是完备的.\par
\textbf{Proof}:设$\{x_n\}$是$X$中的Cauchy列,即$\forall \varepsilon>0$,存在$N_1>0$,使得对任意$m>n>N_1$都有
$$
\|x_m-x_n\|<\varepsilon
$$
下面说明$\{x_n+V\}$是$X/V$中的Cauchy列.事实上
$$
\begin{aligned}
    \|(x_m+V)-(x_n+V)\|=\|(x_m-x_n)+V\| &=\inf{\|(x_m-x_n)+v\|,v\in V}\\
                                        &\leqslant \|x_m-x_n\| 
\end{aligned}
$$
(关于不等号是因为下确界是对所有$v\in V$取的,我们直接取$v=0$)这就表明$\{x_n+V\}$是$X/V$中的Cauchy列,不妨设其收敛到$x_0+V$.这事实上就是在说
$$
\|(x_n-x_0)+V\|=\inf\{\|x_n-x_0+v\|,v\in V\}\to 0\quad n\to \infty
$$
因此根据下确界性质,我们总可以找到$\{v_n\}\subset V$使得对上述$\varepsilon>0$,存在$N_2>0$使得当$n>N_2$时有
$$
\|x_n-x_0+v_n\|\leqslant \frac{\varepsilon}{2}
$$
若记$y_n=x_n+v_n$,上面其实就是在说$\{y_n\}$是$X$中的收敛序列,自然也是$X$中的Cauchy列.注意到$v_n=y_n-x_n$,这也表明$\{v_n\}$是$V$中的Cauchy列,而$V$是完备的,所以存在$v_0\in V$,使得对上述$\varepsilon>0$,存在$N_3>0$,使得对于任意$n>N_3$都有$\|v_n-v_0\|<\frac{\varepsilon}{2}$.\par
取$N=\max\{N_1,N_2,N_3\}$,则当$n>N$时有
$$
\|x_n-x_0+v_0\|\leqslant \|y_n-x_0\|+\|v_n-v_0\|<\varepsilon
$$
显然$(x_0-v_0)\in X$.这就说明了$X$是完备的.\par
在泛函分析中,一个重要的部分就是计算算子范数.一般情况下分为两步:第一步总是通过各种放缩手段(这一过程为了凑右侧的$\|x\|$)建立不等式$\|Tx\|\leqslant M\|x\|$,得到$\|T\|\leqslant M$;第二步主要证明相反的不等式$\|T\|\geqslant M$,最常见的是取特殊的$x_0$满足$\|x_0\|=1$且$\|Tx_0\|=M$,如果看教材上的例子会发现,还可以构造序列$\{x_n\}$满足$\|x_n\|=1$且$\|Tx_n\|\to M$.从而利用算子范数定义成立
$$
\|T\|\geqslant \|Tx_0\|=M\text{或}\|T\|\geqslant \sup_{n}\|Tx_n\|=M
$$
再次强调,第二步的本质上是要证明反向不等式$\|T\|\geqslant M$.\par
\textbf{Problem 8}$^\clubsuit $:求Hardy算子
$$
(Tf)(x)=\frac{1}{x}\int_0^x f(y)\mathrm{d}y
$$
在$L^p(0,\infty)\,(1<p<\infty)$上的范数.\par
\textbf{Proof}:直接计算得到
$$
\begin{aligned}
    \|Tf\|_{L^p} &= \left(\int_0^\infty\left|\frac{1}{x}\int_0^x f(y)\mathrm{d}y\right|^p\right)^{1/p}\\
                 &= \left(\int_0^\infty\left|\frac{1}{x}\int_0^1 f(ux)\mathrm{d}u\right|^p\right)^{1/p}\\
                 &\leqslant \int_0^1\left(\int_0^\infty |f(ux)|^p\mathrm{d}x\right)^{1/p}\mathrm{d}u\\
                 &= \int_0^1\left(\int_0^\infty |f(y)|^p\frac{1}{u}\mathrm{d}x\right)^{1/p}\mathrm{d}u\\
                 &= \left(\int_0^1u^{-1/p}\mathrm{d}u\right)\left(\int_0^\infty|f(y)|^p\mathrm{d}y\right)^{1/p}=\frac{p}{p-1}\|f\|_{L^p}
\end{aligned}
$$
从而$\|T\|\leqslant \frac{p}{p-1}$.\par
对$\forall \varepsilon\in (0,1)$,取
$$
f_0(x)=(1-\varepsilon+x)^{(\varepsilon-2)/p}
$$
则$\|f_0\|_{L^p}=\left(\int_0^{\infty}(1-\varepsilon+x)^{(\varepsilon-2)}\mathrm{d}x\right)^{1/p}=(1-\varepsilon)^{\varepsilon-2}$.因此
$$
\begin{aligned}
    \|Tf_0\|_{L^p} &= \left(\int_0^{\infty}\left|\frac{1}{x}\int_0^x (1-\varepsilon+y)^{(\varepsilon-2)/p}\mathrm{d}y\right|^p\mathrm{d}x\right)^{1/p}\\
                   &= \left(\int_0^{\infty}\left|\int_0^1 (1-\varepsilon+ux)^{(\varepsilon-2)/p}\mathrm{d}u\right|^p\mathrm{d}x\right)^{1/p}\\
                   &\geqslant \left(\int_0^{\infty}\left|\int_0^1 \left((\sqrt{1-\varepsilon}+x)(\sqrt{1-\varepsilon}+u)\right)^{(\varepsilon-2)/p}\mathrm{d}u\right|^p\mathrm{d}x\right)^{1/p}\\
                   &= \left(\int_0^1(\sqrt{1-\varepsilon}+u)^{(\varepsilon-2)/p}\mathrm{d}u \right)\left(\int_0^{\infty}(\sqrt{1-\varepsilon}+x)^{(\varepsilon-2)}\mathrm{d}x \right)^{1/p}\\
                   &= C_p(\varepsilon)(1-\varepsilon)^{(\varepsilon-3)/2}
\end{aligned}
$$
其中$C_p(\varepsilon)=\left(\int_0^1(\sqrt{1-\varepsilon}+u)^{(\varepsilon-2)/p}\mathrm{d}u \right)$.于是
$$
\|T\|\geqslant \frac{\|Tf_0\|_{L^p}}{\|f_0\|_{L^p}}\geqslant \frac{C_p(\varepsilon)(1-\varepsilon)^{(\varepsilon-3)/2}}{(1-\varepsilon)^{\varepsilon-2}}=C_p(\varepsilon)(1-\varepsilon)^{(1-\varepsilon)/2}
$$
令$\varepsilon\to 1$,得到
$$
\|T\|\geqslant \int_0^1u^{-1/p}\mathrm{d}u=\frac{p}{p-1}
$$
故$\|T\|=\frac{p}{p-1}$.事实上我们得到了著名的Hardy不等式
$$
\left(\int_0^\infty\left|\frac{1}{x}\int_0^x f(y)\mathrm{d}y\right|^p\right)^{1/p}\leqslant \frac{p}{p-1}\left(\int_0^\infty|f(y)|^p\mathrm{d}y\right)^{1/p}
$$
式中$\|T\|=\frac{p}{p-1}$是最佳常数.\par
\textbf{Problem 9}$^\clubsuit $:设$(X,d)$是具有以下性质的度量空间:给定$X$的任一非空闭子集序列$\{A_n\}_{n=0}^{\infty}$,它满足对所有的$n\geqslant 0$,$A_n\supset A_{n+1}$且$\lim_{n\to\infty}\operatorname{diam}A_n=0$,则其交集$\cap_{n\geqslant 0}A_n$定为非空.证明:$(X,d)$是完备的.\par
\textbf{Proof}:设$\{x_n\}$为$X$中的Cauchy列,则对任意$\varepsilon>0$,都存在$N>0$,使得$\forall m>n\geqslant N$都有
$$
d(x_m,x_n)<\varepsilon
$$
因此我们可以设当$n\geqslant N$时有$d(x_{n+1},x_n)<\frac{\varepsilon}{2^{n+1}}$.这表明$x_{n+1}\in B(x_n,\frac{\varepsilon}{2^{n+1}})\subset \overline{B(x_n,\frac{\varepsilon}{2^n})}$.\par
因此,对于任意$k\in \mathbb{N}$,取$\varepsilon<2^{-k}$,则存在$N_k>0$,使得对任意$m>n\geqslant N_k$都有$d(x_m,x_n)<2^{-k}$.\par
任取$n_1\geqslant N_1$,选$n_2>\max\{n_1,N_2\}$,则$d(x_{n_2},x_{n_1})<2^{-2}$;\par
选$n_3>\max\{n_2,N_3\}$,则$d(x_{n_3},x_{n_2})<2^{-3}$;\par
如此进行下去,我们得到$\{x_n\}$的子列$\{x_{n_k}\}$满足
$$
d(x_{n_{k+1}},x_{n_{k}})<\frac{1}{2^k}
$$
我们记$A_k=\overline{B(x_{n_k}),\frac{1}{2^k}}$,显然,对$k\geqslant 1$有$A_k\supset A_{k+1}$且$\operatorname{diam}(A_k)\to 0(k\to\infty)$.\par
根据条件,存在$x_0\in X$使得$\cap_{k\geqslant 1}A_k=\{x_0\}$.这表明$x_{n_k}\to x_0(k\to\infty)$.所以$(X,d)$完备.\par
\textbf{Problem 10}$^\clubsuit $:设$T$是$X$的线性子空间$D(T)$到$Y$中的线性算子,则逆算子$T^{-1}:R(T)\to D(T)$存在的充要条件是$N(T)=\{0\}$.\par
\textbf{Proof}:充分性.设$N(T)=\{0\}$.即$\forall x_1,x_2\in D(T),Tx_1=Tx_2\Rightarrow T(x_1-x_2)=0$,故$x_1-x_2\in N(T)=\{0\}$,即$x_1=x_2$.于是$T:D(T)\to R(T)$是双射,于是存在逆算子$T^{-1}:R(T)\to D(T)$.\par
必要性.设$T^{-1}$存在,则$\forall x_1,x_2\in D(T),Tx_1=Tx_2\Rightarrow x_1=x_2$.这表明,$\forall x\in D(T),Tx=0\Rightarrow x=0$,即$N(T)=\{0\}$.\par
\textbf{Problem 11}$^\clubsuit $:设$X,Y$为赋范线性空间,满射$T:X\to Y$是线性算子.若存在正的常数$M$使得
$$
\|Tx\|\geqslant M\|x\|,\forall x\in X
$$
则逆算子$T^{-1}:Y\to X$存在且有界.\par
\textbf{Proof}:先证明$T^{-1}$的存在性,只要证明$N(T)=\{0\}$即可.显然$\{0\}\subset N(T)$,下面证明$N(T)\subset\{0\}$.$\forall x\in N(T)\Rightarrow Tx=0$,从而
$$
M\|x\|\leqslant \|Tx\|=0\Rightarrow x=0
$$
即$N(T)=\{0\}$.\par
由$R(T)=Y$,故对任意$y\in Y$,存在$x\in X$使$y=Tx,x=T^{-1}y$,所以
$$
\|T^{-1}y\|=\|x\|\leqslant \frac{1}{M}\|Tx\|=\frac{1}{M}\|y\|
$$
即$\|T^{-1}\|\leqslant \frac{1}{M}$,这就证明了$T^{-1}$是有界的.\par
\textbf{Problem 12}$^\clubsuit $:设$X$是赋范线性空间,$f$是$X$上的线性泛函,则$f$有界的充要条件是$f$的零空间$N(f)$是$X$中闭线性子空间.\par
\textbf{Proof}:必要性.$\forall\{x_n\}\subset N(f)$满足$x_n\to x(n\to\infty)$.下面证明$x\in N(f)$.由$f$有界,则连续,于是
$$
Tx=T(\lim_{n\to\infty}x_n)=\lim_{n\to\infty}Tx_n=0
$$
即$x\in N(f)$,从而$N(f)$闭.\par
接下来是充分性,设$N(f)$是$X$的闭子空间.若$f$无界,则$\forall n$,存在$x_n\in X$使得
$$
|f(x_n)|> n\|x_n\|
$$
令$y_n=\frac{x_n}{f(x_n)}-\frac{x_1}{f(x_1)}$,则$f(y_n)=0$,故$y_n\in N(f)$.记$z=-\frac{x_1}{f(x_1)}$,则
$$
\|y_n-z\|=\left\|\frac{x_n}{|f(x_n)|}\right\|<\frac{1}{n}\to 0\quad (n\to \infty)
$$
这说明$y_n\to z$.注意到$N(f)$是闭的,所以$z\in N(f)$,那么应该有$f(z)=0$.但$f(z)=-1\neq 0$,矛盾.因此$f$有界.\par
\textbf{Problem 13}$^\clubsuit $:设$X$是赋范线性空间,$f$是$X$上的线性泛函,则非零线性泛函$f$有界的充要条件是$f$的零空间$N(f)$不在$X$中稠密.\par
\textbf{Proof}:因为$f\neq 0$,所以$N(f)\neq X$.又$f$有界$\Leftrightarrow $$N(f)$闭,于是$\overline{N(f)}=N(f)\neq X$.\par
\textbf{Problem 14}$^\clubsuit $:对任意$f\in L[a,b]$,定义积分算子$T$;
$$
(Tf)(x)=\int_a^xf(t)\mathrm{d}t
$$
\qquad (1).将$T$看成$L[a,b]\to C[a,b]$的算子时,$\|T\|=1$.\par
(2).将$T$看成$L[a,b]\to L[a,b]$的算子时,$\|T\|=b-a$.\par
\textbf{Proof}:(1).利用范数的定义
$$
\|Tf\|=\max_{[a,b]}|(Tf)(x)|=\max_{[a,b]}\left|\int_a^xf(t)\mathrm{d}t\right|\leqslant \max_{[a,b]}\int_a^x|f(t)|\mathrm{d}t\leqslant \int_a^b|f(t)|\mathrm{d}t=\|f\|_{L^1}
$$
这表明$\|T\|\leqslant 1$.\par
另一方面,取$f_0(t)=\frac{1}{b-a},t\in [a,b]$.则$\|f_0\|_{L^1}=1$,且
$$
\|T\|=\sup_{\|f\|_{L^1}=1}\|Tf\|\geqslant \|Tf_0\|=\max_{[a,b]}\int_a^x\frac{1}{b-a}\mathrm{d}t=1
$$
从而$\|T\|=1$.\par
(2).依然利用范数定义
$$
\|Tf\|=\int_a^b\left|\int_a^xf(t)\mathrm{d}t\right|\mathrm{d}x\leqslant \int_a^b\left(\int_a^x|f(t)|\mathrm{d}t\right)\mathrm{d}x\leqslant \int_a^b\left(\int_a^b|f(t)|\mathrm{d}t\right)\mathrm{d}x=(b-a)\|f\|
$$
这表明$\|T\|\leqslant b-a$.\par
对任何满足$a+\frac{1}{n}<b$的自然数$n$,令
$$
f_n(t)=
\begin{cases}
    n,&a\leqslant t\leqslant a+\frac{1}{n}\\[5pt]
    0,&a+\frac{1}{n}<t\leqslant b
\end{cases}
$$
则$\|f_n\|=1$.下面将证明$\|Tf_n\|\to b-a(n\to\infty)$.事实上
$$
\begin{aligned}
    \|Tf_n\| &= \int_a^b\left|\int_a^xf_n(t)\mathrm{d}t\right|\mathrm{d}x\\
             &= \int_{a}^{a+1/n}n(x-a)\mathrm{d}x+\int_{a+1/n}^{b}\left|\int_{a}^{a+1/n}n\mathrm{d}t+\int_{a+1/n}^{b}0\mathrm{d}t\right|\mathrm{d}x\\
             &= b-a-\frac{1}{2n}
\end{aligned}
$$
因此$\|T\|\geqslant \sup\{\|Tf_n\|:n>\frac{1}{b-a}\}=b-a$.从而$\|T\|=b-a$.\par
关于算子空间,我们尤其要理解算子的各种收敛定义.如果我们设$X,Y$为数域$\mathbb{K}$上的赋范线性空间,$T,T_n\in B(X,Y)$.那么\par
(1).$T_n\Rightarrow T$(依算子范数或一致收敛)$\Longleftrightarrow \lim_{n\to\infty}\|T_n-T\|=0$;\par
(2).$T_n \xrightarrow{s} T$(强收敛或点态收敛)$\Longleftrightarrow \forall x\in X,\lim_{n\to\infty}\|T_nx-Tx\|=0$;\par
(3).$T_n \xrightarrow{w} T$(弱收敛)$\Longleftrightarrow \forall x\in X,f\in Y',\lim_{n\to\infty}\|f(T_nx)-f(Tx)\|=0$;\par
需要注意的是,当$Y=\mathbb{K}$的时候,得到泛函的收敛,此时称呼不一样.如果我们设$X$为赋范线性空间,$f,f_n\in X'$.那么\par
(1).$f_n\xrightarrow{s} f$(强收敛)$\Longleftrightarrow \lim_{n\to\infty}\|f_n-f\|=0$;\par
(2).$f_n \xrightarrow{w} f$(弱收敛)$\Longleftrightarrow \forall g\in X'',\lim_{n\to\infty}\|g(f_n)-g(f)\|=0$;\par
(3).$f_n \xrightarrow{w^*} f$(弱$*$收敛)$\Longleftrightarrow \forall x\in X,\lim_{n\to\infty}\|f_n(x)-f(x)\|=0$;\par
\textbf{Problem 15}$^\clubsuit $:设$X$是Banach空间,$Y$是赋范线性空间,$T_n\in B(X,Y)$,若$T_n\xrightarrow{s} T$,则$T\in B(X,Y)$且
$$
\|T\|\leqslant \liminf_{n\to\infty}\|T_n\|
$$
\qquad \textbf{Proof}:由$T_n\in B(X,Y)$且$T_n\xrightarrow{s} T$,因此$\{T_n\}$是点态有界的,由一致有界定理知$\sup_{n}\|T_n\|<\infty$,则存在$M>0$使$\forall n$有$\|T_n\|\leqslant M$.由三角不等式得
$$
\left|\|T_nx\|-\|Tx\|\right|\leqslant \|T_nx-Tx\|\to 0\quad n\to \infty
$$
这说明$\lim_{n\to \infty}\|T_nx\|=\|Tx\|$.又$\|T_nx\|\leqslant\|T_n\|\|x\|$.令$n\to\infty$得到
$$
\|Tx\|=\lim_{n\to \infty}\|T_nx\|\leqslant \left(\liminf_{n\to\infty}\|T_n\|\right)\|x\|
$$
于是$\|T\|=\sup_{x\neq 0}\frac{\|Tx\|}{\|x\|}\leqslant \liminf_{n\to\infty}\|T_n\|$.\par
\textbf{Problem 16}$^\clubsuit $:设$(X,d)$为度量空间,$E$是$X$的子集.$f\in C(E)$的充要条件是对$\forall \alpha\in \mathbb{R}$,$f$的水平集$\{f\geqslant \alpha\}$或者$\{f\leqslant\alpha\}$均为$E$中闭集,或等价地$\{f>\alpha\}$和$\{f<\alpha\}$为$E$中开集.\par
\textbf{Proof}:必要性.若$f\in C(E)$.记$F=\{f\geqslant\alpha\}$.若$F$的导集$F'=\emptyset$,此时当然有$F'\subset F$,因此$F$是闭集.下设$F'\neq\emptyset$.取$x_0\in F'$,那么存在$\{x_n\}\subset F$使得$x_n\to x_0(n\to\infty)$.因为$f(x_n)\geqslant\alpha$,根据$f\in C(E)$,得到
$f(x_0)=\lim_{n\to\infty}f(x_n)\geqslant\alpha$,这就表明$x_0\in F$.因此$F'\subset F$,故$F$是闭集,从而有$F^c=\{f<\alpha\}$为开集.\par
充分性.设对$\forall \alpha_1,\alpha_2\in\mathbb{R},G_1=\{f<\alpha_1\},G_2=\{f<\alpha_2\}$为开集.特别的,对$\forall x_0\in E,\forall\varepsilon>0$,取$\alpha_1=f(x_0)+\varepsilon,\alpha_2=f(x_0)-\varepsilon$,则$x_0\in G_1\cap G_2$.由$G_1\cap G_2$为开集,存在$\delta>0$,使得$B(x_0,\delta)\subset G_1\cap G_2$,这说明$\forall x\in B(x_0,\delta)$成立$x\in G_1\cap G_2$,即
$$
f(x_0)-\varepsilon<f(x)<f(x_0)+\varepsilon
$$
这表明$f$在$x_0$处连续.由$x_0$的任意性,$f\in C(E)$.\par
\textbf{Problem 17}$^\clubsuit $:设$\{x_n\}$是Banach空间$X$的点列,$x\in X$.若$x_n\xrightarrow{w} x$且$\sup_{n}\|x_n\|<\infty$,则$\|x\|\leqslant\liminf_{n\to\infty}\|x_n\|$.\par
\textbf{Proof}:由题意,$\forall f\in X'$,有$\lim_{n\to\infty}|f(x_n)-f(x)|=0$.由三角不等式
$$
\left||f(x_n)|-|f(x)|\right|\leqslant |f(x_n)-f(x)|\to 0,\quad n\to\infty
$$
这表明$|f(x_n)|\to |f(x)|(n\to\infty)$.又因为
$$
|f(x_n)|\leqslant \|f\|\|x_n\|
$$
令$n\to\infty$得到$|f(x)|\leqslant \|f\|(\liminf_{n\to\infty}\|x_n\|)$.所以
$$
\|x\|=\sup\{|f(x)|:f\in X',\|f\|=1\}\leqslant \liminf_{n\to\infty}\|x_n\|
$$
\qquad \textbf{Problem 18}$^\clubsuit $:$G$是赋范线性空间$X$的子空间,$x_0\in X$.则$x_0\in\overline{G}$当且仅当对于任意满足$f|_{G}=0$的$f\in X'$必有$f(x_0)=0$.\par
\textbf{Proof}:只证充分性.不妨设$x_0\notin \overline{G}$.定义映射$\phi:G+\mathbb{K}x_0\to \mathbb{K}$如下
$$
\phi(g+\alpha x_0)=\alpha,\forall g\in G,\alpha\in\mathbb{K}
$$
那么我们定义了一个$G+\mathbb{K}x_0$上的线性泛函.显然$N(\phi)=G$,$\phi(x_0)=1$.\par
因为$x_0\notin \overline{G}$,所以$\overline{N(\phi)}=\overline{G}\neq \overline{G+\mathbb{K}x_0}$,这说明$\phi$是连续的.由Hahn-Banach延拓定理可知,存在$f_0\in X'$满足在$G+\mathbb{K}x_0$上$f_0=\phi$.容易看出$f_0(x_0)=\phi(x_0)=1$且$f_0|_{G}=0$.\par
\textbf{Problem 19}$^\clubsuit $:$(X,\|\cdot\|)$是无穷维赋范线性空间,则$X'$也是无穷维的.\par
\textbf{Proof}:设$\operatorname{dim}X=\infty$,所以$X$中存在线性无关列$\{e_k\}$.令$A_n=\operatorname{span}\{e_1,\cdots,e_n\}$,这说明$\operatorname{dim}A_n<\infty$,而任意无穷维赋范线性空间的有限维子空间都是闭的,故$A_n$是$X$的真闭子空间,则有$e_{n+1}\in X-A_n$,所以存在$f_n\in X'$满足对$\forall x\in A_n,f_n(x)=0$且$f_n(e_{n+1})=d(e_{n+1},A_n)=d_n>0$.\par
下面证明$\{f_n\}$是$X'$中的线性无关列.设$\forall c_k(1\leqslant k\leqslant m),\forall m\in\mathbb{N}$满足$\sum_{k=1}^{m}c_kf_k=0$.对$\forall e_k\in A_m,1\leqslant k\leqslant m$,有
$$
c_1d_1=c_1f_1(e_2)=\sum_{k=1}^{m}c_kf_k(e_2)=0
$$
而$d_1>0$,所以$c_1=0$且$\sum_{k=2}^{m}c_kf_k(e_2)$.同理
$$
c_2d_2=c_2f_2(e_3)=\sum_{k=2}^{m}c_kf_k(e_3)=0
$$
得$c_2=0$.用同样的方法可以得出$c_k=0,1\leqslant k\leqslant m$.这就表明$\{f_1,f_2,\cdots,f_m\}$是线性无关的序列,再由$m$的任意性就可以得到$\{f_n\}$是$X'$中的线性无关列.\par
\textbf{Problem 20}$^\clubsuit $:设$X,Y$为赋范线性空间,$T:X\to Y$为线性算子.若$T$是闭算子且逆算子$T^{-1}:Y\to X$存在,证明:$T^{-1}$也是闭算子.\par
\textbf{Proof}:根据闭算子的定义:即证明任意$\{y_n\}\subset Y$满足$y_n\to y_0,T^{-1}y_n\to x_0$,那么$y_0\in Y$且$T^{-1}y_0=x_0$.\par
由题意可以看出$T$是双射,于是对每个$y_n\in Y$都存在$x_n\in X$使得$Tx_n=y_n$,即$x_n=T^{-1}y_n$,这说明$x_n\to x_0$,且$Tx_n=y_n\to y_0$.因为$T$是闭算子,由$x_n\to x_0,Tx_n\to y_0$可以得到$x_0\in X$且$Tx_0=y_0$,于是$T^{-1}y_0=x_0$.\par
再用$T$是双射且$Tx_0=y_0$,显然$y_0\in Y$.因此$T^{-1}$是闭算子.\par
\textbf{Problem 21}$^\clubsuit $:设$\{a_n\}$是复数列,使得对于任何收敛于零的复数列$\{x_n\}$,级数$\sum_{n=1}^{\infty}a_nx_n$都收敛.证明$\sum_{n=1}^{\infty}|a_n|<\infty$.\par
\textbf{Proof}:我们令$c_0$表示所有收敛于零的序列全体,则其显然是完备的.在$c_0$上定义线性泛函如下
$$
f_n(x)=\sum_{k=1}^{n}a_kx_k,\forall x=\{x_1,\cdots,x_n,\cdots\}\in c_0
$$
于是
$$
|f_n(x)|\leqslant \sum_{k=1}^{n}|a_k||x_k|\leqslant \left(\sum_{k=1}^{n}|a_k|\right)\|x\|
$$
所以$\|f_n\|\leqslant \sum_{k=1}^{n}|a_k|$.取$x_0=\{\sgn a_1,\cdots,\sgn a_n,0,\cdots \}$,则$x_0\in c_0,\|x\|=1$,且
$$
\|f_n\|\geqslant |f_n(x_0)|=\sum_{k=1}^{n}|a_k|
$$
故$\|f_n\|=\sum_{k=1}^{n}|a_k|$.这说明$f_n\in (c_0)'=l^1$.而由假设$\lim_{n\to\infty}f_n=\sum_{n=1}^{\infty}a_nx_n$收敛,从而$\sup_{n}|f_n(x)|<\infty$.由一致有界定理知$\sup_{n}\|f_n\|<\infty$,即$\sum_{n=1}^{\infty}|a_n|=\sup_{n}\|f_n\|<\infty$.\par
\textbf{Problem 22}$^\clubsuit $:设$M$是Hilbert空间$H$的子集,求证:$M^{\perp\perp}=\overline{\operatorname{span}M}$.\par
\textbf{Proof}:很容易观察到下面的关系
$$
\begin{aligned}
    x\in M^{\perp} \Leftrightarrow x\perp M \Leftrightarrow x\perp \operatorname{span}M &\Leftrightarrow x\perp \overline{\operatorname{span}M}\\
                                                                                        &\Leftrightarrow x\in \overline{\operatorname{span}M}^{\perp}
\end{aligned}
$$
这说明$M^{\perp}=\overline{\operatorname{span}M}^{\perp}$,如果要证$M^{\perp\perp}=\overline{\operatorname{span}M}$,即证明$\overline{\operatorname{span}M}=\overline{\operatorname{span}M}^{\perp\perp}$.问题归结为证明:若$M$是闭的,则$M=M^{\perp\perp}$.\par
我么已知$M\subset M^{\perp\perp}$,下面只需要证明$M^{\perp\perp}\subset M$.任取$x\in M^{\perp\perp}$,由正交分解定理
$$
x=x_0+y,x_0\in M,y\in M^{\perp}
$$
又$y\in M^{\perp},x\in M^{\perp\perp}$,因此$(x,y)=0$,从而$0=(x,y)=(x_0,y)+(y,y)$.又因为$(x_0,y)=0$,这说明$y=0$.于是$x=x_0\in M$.综上我们证明了$M=M^{\perp\perp}$.\par
\textbf{Problem 23}$^\clubsuit $:设$X,Y,Z$是Banach空间,若$T_1\in B(X,Z),T_2\in B(Y,Z)$,且对$\forall x\in X$算子方程$T_1x=T_2y$都有唯一解$y=Tx$.证明$T\in B(X,Y)$.\par
\textbf{Proof}:容易验证$T:X\to Y$是线性算子.注意到$X,Y$都是Banach空间,我们只需要证明$T$是闭算子,再根据闭图像定理就能得到$T$有界.\par
任取$\{x_n\}\subset X$,满足$x_n\to x_0,Tx_n\to y_0(n\to\infty)$.下面证明$x_0\in X,Tx_0=y_0$.因为$X$为Banach空间,故一定有$x_0\in X$.记$y_n=Tx_n$,由假设条件$T_1x_n=T_2y_n$与$T_1,T_2$的连续性知$T_1x_0=T_2y_0$,且这个方程有唯一解$y_0=Tx_0$.这就证明了$T$为闭算子.\par
\textbf{Problem 24}$^\clubsuit $:设$X,Y$为Banach空间,$T\in B(X,Y)$.若$T$是双射,证明存在正的常数$a,b$使得
$$
a\|x\|\leqslant \|Tx\|\leqslant b\|x\|,\forall x\in X
$$
\qquad \textbf{Proof}:由逆算子定理可知$T^{-1}\in B(Y,X)$,从而
$$
1=\|TT^{-1}\|\leqslant \|T\|\|T^{-1}\|\Rightarrow \|T\|>0,\|T^{-1}\|>0
$$
令$a=\|T^{-1}\|^{-1},b=\|T\|$,则$a>0,b>0$.对任意$x\in X$,记$Tx=y$,则$x=T^{-1}y$,这说明
$$
\|x\|=\|T^{-1}(Tx)\|\leqslant \|T^{-1}\|\|Tx\|\Rightarrow a\|x\|\leqslant \|Tx\|
$$
左边不等式成立.又$\|Tx\|\leqslant \|T\|\|x\|=b\|x\|$,右边不等式成立.\par
\textbf{Problem 25}$^\clubsuit $:设$F$是赋范线性空间$(X,\|\cdot\|)$的闭子空间.若$x_n\in F$且$x_n\xrightarrow{w}x_0$.求证$x_0\in F$.\par
\textbf{Proof}:设$x_0\notin F$,而$F$为$X$的闭子空间,则$d(x_0,F)=d>0$.于是,存在$f\in X'$满足$f|_{F}=0$且$f(x_0)=d>0$.又$x_n\in F$,则$f(x_n)=0$.根据题设$x_n\xrightarrow{w} x_0$,则由定义,$f(x_0)=\lim_{n\to \infty}f(x_n)=0$,但这和$f(x_0)=d>0$矛盾.所以$x_0\in F$.\par
\textbf{Problem 26}$^\clubsuit $:设$X$是复Hilbert空间,$\{T_n\}$是$X$上自伴随算子,若存在算子$T:X\to X$使得
$$
\lim_{n\to\infty}(T_nx,y)=(Tx,y),\forall x,y\in X
$$
则$T$也是自伴随的.\par
\textbf{Proof}:由$T_n$是自伴随的,故$T_n^{*}=T_n$.根据定义$(T_nx,y)=(x,T_n^{*}y)=(x,T_ny)$,从而
$$
\begin{aligned}
    (Tx,y)=\lim_{n\to\infty}(Tx_n,y) &= \lim_{n\to\infty}(x,T_ny)\\
                                     &= \lim_{n\to\infty}\overline{(T_ny,x)}=\overline{(Ty,x)}=(x,Ty)
\end{aligned}
$$
由Hilbert空间伴随算子的存在唯一性知$(Tx,y)=(x,T^{*}y)$,于是对任意$x,y\in X,(x,(T-T^{*})y)=0$,这就证明了$T=T^{*}$.\par
\textbf{Problem 27}$^\clubsuit $:设$(X,d)$为可分的度量空间,$E$是$X$的子空间,则$E$也是可分的.\par
\textbf{Proof}:由$X$可分,故$X$中存在可数的稠密子集$D$.记$D=\{x_1,x_2\cdots,x_n,\cdots\}$.令$E_k=B(x_k,\frac{1}{k})\cap E$.任取$y_k\in E_k$,令$A=\{y_k\}$,则$A$是$E$中可数集.\par
下面证明$E\subset\overline{A}$.对$\forall x_0\in E\subset X=\overline{D}$,那么存在$x_k\in D$使得$d(x_k,x_0)\to 0$.又$y_k\in E_k=B(x_k,\frac{1}{k})\cap E$,因此$d(y_k,x_k)<\frac{1}{k}\to 0$.于是
$$
d(y_k,x_0)\leqslant d(y_k,x_k)+d(x_k,x_0)\to 0,quad k\to\infty
$$
即$y_k\to x_0$.这就说明$x_0\in\overline{A}$.从而$E=\overline{A}$,所以$E$可分.\par
\textbf{Problem 27}$^\clubsuit $:设$H$为Hilbert空间,$A\in B(H,H)$.并且存在$m>0$使得
$$
|(Ax,x)|\geqslant m\|x\|^2
$$
求证$A^{-1}\in B(H,H)$.\par
\textbf{Proof}:由条件,$\forall x\in H$,因为
$$
m\|x\|^2\leqslant |(Ax,x)|\leqslant \|Ax\|\|x\|
$$
则在$x\neq 0$时,$\|Ax\|\geqslant \|x\|$,事实上,这一不等式对$x=0$的情况也成立.由此可以得到$A$是单射.\par
又$\forall x\in (R(A))^{\perp}$,则$0=|(Ax,x)|\geqslant m\|x\|^2\Rightarrow x=0$,即$(R(A))^{\perp}=\{0\}$,这说明$R(A)$在$H$中稠密,即$\overline{R(A)}=H$.要证$A$是满射,即$R(A)=H$,问题可转化为证明$R(A)$是闭集.\par
设$\{y_n\}$是$R(A)$中的收敛序列,所以是Cauchy列.设$Ax_n=y_n$,则根据题设
$$
\|x_m-x_n\|\leqslant \frac{1}{m}\|Ax_m-Ax_n\|=\frac{1}{m}\|y_m-y_n\|
$$
由此可见,$\{x_n\}$也是Cauchy列.由$H$是Hilbert空间,从而$\{x_n\}$收敛.设$x_n\to x_0(n\to\infty)$,于是
$$
y_n=Ax_n\to Ax_0\in R(A)
$$
从而$R(A)$是闭的,即有$R(A)=\overline{R(A)}=H$,这就证明了$A$是满射.由逆算子定理可知,$A^{-1}\in B(H,H)$.\par
\textbf{Problem 28}$^\clubsuit $:设$X$是赋范线性空间,$\{x_{\alpha}:\alpha\in \varLambda \}\subset X$.若对$\forall f\in X'$有$\sup_{\alpha\in \varLambda }\{|f(x_\alpha)|\}<\infty$.证明$\sup_{\alpha\in \varLambda }\{\|x_\alpha\|\}<\infty$.\par
\textbf{Proof}:定义算子
$$
\begin{aligned}
    T_{\alpha}:X' &\to \mathbb{R}\\
               f  &\mapsto f(x_{\alpha})
\end{aligned}
$$
则$T_{\alpha}$是线性算子,且对固定的$\alpha$有
$$
\|T_{\alpha}(f)\|=\|f(x_{\alpha})\|\leqslant\|f\|\|x_{\alpha}\|
$$
这说明$T_{\alpha}$是有界线性算子.又
$$
\sup_{\alpha\in \varLambda }\{\|T_{\alpha}(f)\|\}=\sup_{\alpha\in \varLambda }\{|f(x_{\alpha})|\}<\infty,\forall f\in X'
$$
且$X'$完备,由一致有界定理知$\sup_{\alpha\in \varLambda }\{\|T_{\alpha}\|\}<\infty$,因为
$$
\|T_{\alpha}\|=\sup_{\|f\|=1}|T_{\alpha}(f)|=\sup_{\|f\|=1}|f(x_{\alpha})|=\|x_{\alpha}\|
$$
所以$\sup_{\alpha\in \varLambda }\{\|x_\alpha\|\}<\infty$.\par
\textbf{Problem 29}$^\clubsuit $:设$X$为Hilbert空间,$M$是$X$的子空间,则$M$在$X$稠密当且仅当$M^{\perp}=\{0\}$.\par
\textbf{Proof}:若$M$在$X$中稠密,则$\forall x\in M^{\perp}$,存在$\{x_n\}\subset M$使$x_n\to x(n\to\infty)$.故$(x_n,x)\to (x,x)$.又$(x_n,x)=0$,因此$(x,x)=0$,于是$x=0$,这就得到了$M^{\perp}=\{0\}$.\par
反之,设$M^{\perp}=\{0\}$,即$\overline{M}^{\perp}=\{0\}$.对$\forall x\in X$,由正交分解定理,存在$x_0\in \overline{M},y\in \overline{M}^{\perp}$使
$$
x=x_0+y
$$
又$\overline{M}^{\perp}=\{0\}$,因此$x=x_0\in\overline{M}$,即$X\subset \overline{M}$.\par
\textbf{Problem 30}$^\clubsuit $:设$X,Y$为Hilbert空间,$T\in B(X,Y)$,则
$$
\begin{aligned}
    N(T) &= R(T^{*})^{\perp},\, N(T^{*})=R(T)^{\perp}\\
    \overline{R(T^{*})} &= N(T)^{\perp},\, \overline{R(T)}=N(T^{*})^{\perp}
\end{aligned}
$$
\qquad \textbf{Proof}:对$\forall x\in N(f)$,有$Tx=0$对任意$z\in R(T^{*})$,都存在$y\in Y$满足$z=T^{*}y$,因此
$$
(x,z)=(x,T^{*}y)=(Tx,y)=0
$$
这说明$x\in R(T^{*})^{\perp}$.故$N(T) \subset R(T^{*})^{\perp}$.\par
另一方面,$\forall x\in R(T^{*})^{\perp}$,那么对任意$y\in Y$都有$(Tx,y)=(x,T^{*}y)=0$.若取$y=Tx$,则由$(Tx,TX)=0$得$Tx=0$,这说明$x\in N(T)$.故$N(T) \supset R(T^{*})^{\perp}$.\par
由$N(T) = R(T^{*})^{\perp}$,得到$N(T^{*})=R(T^{**})^{\perp}=R(T)^{\perp}$.\par
由$N(T) = R(T^{*})^{\perp}$,得到
$$
\begin{aligned}
    N(T)^{\perp}=R(T^{*})^{\perp\perp}=\overline{R(T^{*})}^{\perp\perp}=\overline{R(T^{*})}
\end{aligned}
$$
\qquad 由$N(T^{*})=R(T)^{\perp}$知
$$
N(T^{*})^{\perp}=R(T)^{\perp\perp}=(\overline{R(T)})^{\perp\perp}=\overline{R(T)}
$$
\qquad \textbf{Problem 31}$^\clubsuit $:设$f\in C[a,b]$,在$[a,b]$内取$n$个点:$a\leqslant x_1<x_2<\cdots<x_n\leqslant b$.令
$$
w(x)=\prod _{k=1}^{n}(x-x_k),l_k(x)=\frac{w(x)}{(x-x_k)w'(x_k)}
$$
则定义如下Lagrange插值算子:
$$
(L_nf)(x)=\sum_{k=1}^{n}f(x_k)l_k(x)
$$
证明:$L_n:C[a,b]\to C[a,b]$是有界线性算子,且$\|L_n\|=\max_{a\leqslant x\leqslant b}\sum_{k=1}^{n}|l_k(x)|$.\par
\textbf{Proof}:记$M=\max_{a\leqslant x\leqslant b}\sum_{k=1}^{n}|l_k(x)|$,则$\|L_n(f)\|=\max_{a\leqslant x\leqslant b}|(L_nf)(x)|\leqslant M\|f\|$.从而有$\|L_n\|\leqslant M$.\par
另一方面,因为$\sum_{k=1}^{n}|l_k(x)|\in C[a,b]$,所以存在$x_0\in [a,b]$使得$\sum_{k=1}^{n}|l_k(x_0)|=M$.选取$f_0\in C[a,b]$,使得$\|f_0\|=1$且$f_0(x_k)=\sgn l_k(x_0),k=1,\cdots,n.|f_0(x)|\leqslant 1$.于是
$$
\begin{aligned}
    \|L_n\| &=\sup_{\|f\|=1}\|L_nf\|\geqslant \|L_n(f_0)\|=\sup_{a\leqslant x\leqslant b}|(L_nf_0)(x)|\\
            &\geqslant |(L_nf_0)(x_0)|=\left|\sum_{k=1}^{n}l_k(x_0)\sgn l_K(x_0)\right|=\sum_{k=1}^{n}|l_k(x_0)|=M
\end{aligned}
$$
\qquad \textbf{Problem 32}$^\clubsuit $:在赋范线性空间中点列的弱收敛无法推出强收敛.\par
\textbf{Proof}:设$X=l^2$,定义点列$\{e_n\}$如下
$$
e_n=(\underbrace{0,0,\cdots,0}_{n-1},1,0,\cdots)
$$
则$\|e_n\|_{l^2}=1$,于是$\{e_n\}$不可能强收敛于0.\par
但对任意$y\in (l^2)'=l^2,y=(\eta_1,\eta_2,\cdots),\sum_{n=1}^{\infty}|\eta_n|^2<\infty$.有
$$
(y,e_n)=\eta_n\to 0,n\to\infty
$$
这说明$\{e_n\}$弱收敛于0.\par
\textbf{Problem 33}$^\clubsuit $:算子的强收敛无法推出一致收敛.\par
\textbf{Proof}:设$X=Y=l^2$,定义算子$T_n:X\to Y$如下
$$
\begin{aligned}
    T_n:X&\to Y\\
    (\xi_1,\xi_2,\cdots)&\mapsto (\underbrace{0,0,\cdots,0}_{n},\xi_{n+1},\xi_{n+2},\cdots)
\end{aligned}
$$
则每个$T_n$都是有界线性算子,且$\|T_n\|\leqslant 1$.此时对任意$x=(\xi_1,\xi_2,\cdots)\in l^2$
$$
\|T_nx-0\|^2=\sum_{k=n+1}^{\infty}|\xi_k|^2\to 0,(n\to\infty)
$$
这说明$T_n$强收敛于0(这个0是零算子).\par
若取$e_{n+1}=(\underbrace{0,0,\cdots,0}_{n},1,0,\cdots)$,那么$T_ne_{n+1}=e_{n+1}$.又因为$\|e_{n+1}\|=1$,容易得到$\|T_n\|=1$,这说明$T_n$不可能一致收敛于0.\par
\textbf{Problem 34}$^\clubsuit $:算子的弱收敛无法推出强收敛.\par
\textbf{Proof}:设$X=Y=l^2$,定义算子$T_n:X\to Y$如下
$$
\begin{aligned}
    T_n:X&\to Y\\
    (\xi_1,\xi_2,\cdots)&\mapsto (\underbrace{0,0,\cdots,0}_{n},\xi_1,\xi_2,\cdots)
\end{aligned}
$$
对任意$x=(\xi_1,\xi_2,\cdots)\in l^2$,有$\|T_nx\|=\|x\|$,于是$T_n$是有界线性算子,且$\|T_n\|=1$.对任意$y\in (l^2)'=l^2,y=(\eta_1,\eta_2,\cdots)$我们有
$$
\begin{aligned}
    (y,T_nx) &=\sum_{j=1}^{\infty}\xi_j\overline{\eta_{n+j}}\\
             &\leqslant \left(\sum_{j=1}^{\infty}|\xi_j|^2\right)^{\frac{1}{2}}\left(\sum_{j=1}^{\infty}|\eta_{n+j}|^2\right)^{\frac{1}{2}}\\
             &\leqslant \|x\|\left(\sum_{k=n+1}^{\infty}|\eta_k|^2\right)^{\frac{1}{2}}\to 0,(n\to\infty)
\end{aligned}
$$
这说明$T_n$弱收敛于0(这个0是零算子).\par
取$e_1=(1,0,\cdots)$,那么当$m\neq n$时,
$$
\|T_ne_1-T_me_1\|=\|e_{n+1}-e_{m+1}\|=\sqrt{2}
$$
这说明$\{T_ne_1\}$不收敛,从而$T_n$不可能强收敛于0.\par
\textbf{Problem 35}$^\clubsuit $:一致有界原理中的完备性条件不能去掉.\par
\textbf{Proof}:设$X$是$l^{\infty}$空间中只有有限多个非零项的序列组成的子空间,定义$T_n:X\to l^{\infty},x=(x_1,x_2,\cdots)\mapsto y=(0,\cdots,0,nx_n,0,\cdots)$,则$T_n\in B(X,l^{\infty})$.\par
对任意$x=(x_1,x_2,\cdots)\in X$,因为$x$的坐标中只有有限多个非零项,所以存在$n_0$,使得$\forall n>n_0$时,$x_n=0$,从而$\|T_nx\|=n|x_n|=0$;对$\forall n\leqslant n_0,\|T_nx\|=n|x_n|\leqslant n_0\|x\|$.于是,对任意$x\in X,\sup_{n}\|T_nx\|<\infty$,即$\{T_n\}$在$X$上点态有界,但$\|T_n\|=n\to\infty(n\to\infty)$.所以$\{T_n\}$不是一致有界的.\par
\textbf{Problem 36}$^\clubsuit $:设$(X,\|\cdot\|)$是赋范线性空间,$X\neq\{0\}$.求证$X$为Banach空间的充要条件是$X$中的单位球面$S=\{x\in X:\|x\|=1\}$是完备的.\par
\textbf{Proof}:只考虑证明充分性.设$\{x_n\}$是$X$中任一Cauchy列,则对任意$\varepsilon>0$,存在$N>0,\forall m,n>N$都有
$$
|\|x_m\|-\|x_n\||\leqslant \|x_m-x_n\|<\varepsilon
$$
这说明$\{\|x_n\|\}$为$\mathbb{R}$中的Cauchy列,不妨设$\{\|x_n\|\}$收敛于$a\in\mathbb{R}$.若$a=0$,则必有$x_n\to 0(n\to\infty)$,显然$0\in X$.下设$a\neq 0$,则存在$N_1>0$使得当$n>N_1$时有$\|x_n\|>\frac{a}{2}$.\par
设$y_n=\frac{x_n}{\|x_n\|}$,则$\|y_n\|=1$.又$\{\|x_n\|\}$是有界序列,即存在$M>0$使得$\|x_n\|\leqslant M$.对$\forall m,n>\max\{N,N_1\}$时
$$
\begin{aligned}
    \|y_m-y_n\| &= \left\|\frac{x_m}{\|x_m\|}-\frac{x_n}{\|x_n\|}\right\|\\
                &=\left\|\frac{\|x_n\|x_m-\|x_m\|x_n}{\|x_m\|\|x_n\|}\right\|\\
                &\leqslant \frac{1}{\|x_m\|\|x_n\|}(\|x_n\|\|x_m-x_n\|+|\|x_n\|-\|x_m\||\|x_n\|)\\
                &<\frac{4M}{a^2}[\|x_m-x_n\|+|\|x_m\|-\|x_n\||]<\frac{8M}{a^2}\varepsilon
\end{aligned}
$$
这说明$\{y_n\}$是$S$中的Cauchy列,由$S$的完备性知,存在$y_0\in S$使得$y_n\to y_0(n\to\infty)$.记$x_0=ay_0$,于是
$$
\begin{aligned}
    \frac{1}{\|x_n\|}\|x_n-x_0\| &= \left\|\frac{x_n}{\|x_n\|}-a\frac{y_0}{\|x_n\|}\right\|\\
                                 &= \left\|y_n-\frac{a}{\|x_n\|}y_0\right\|\to \left\|y_0-\frac{a}{a}y_0\right\|=0,(n\to\infty)
\end{aligned}
$$
即$\{x_n\}$收敛于$x_0=ay_0\in X$,这就证明了$X$是Banach空间.\par